\documentclass[11pt]{article}
% Version française de article.tex ; la version anglaise est article.tex. Compilée avec XeLaTeX :
% cd docs/bridges && xelatex article-fr.tex && xelatex article-fr.tex
\usepackage{polyglossia}\setmainlanguage{french}\setotherlanguage{english}
\newcommand{\thmname}{Théorème}\newcommand{\corname}{Corollaire}\newcommand{\propname}{Proposition}
\newcommand{\lemname}{Lemme}\newcommand{\bridgename}{Pont}\newcommand{\defname}{Définition}
\newcommand{\exname}{Exemple}\newcommand{\exoname}{Exercice}\newcommand{\remname}{Remarque}
% Shared preamble of docs/bridges/article.tex (English) and article-fr.tex (French).
% Built with XeLaTeX: cd docs/bridges && xelatex article.tex && xelatex article.tex
\usepackage{fontspec}
\IfFontExistsTF{Charter}{\setmainfont{Charter}}{\IfFontExistsTF{XCharter}{\setmainfont{XCharter}}{}}
\setmonofont{DejaVu Sans Mono}[Scale=0.8]
\usepackage{microtype}
\usepackage[a4paper,margin=27mm]{geometry}
\usepackage{setspace}\setstretch{1.1}
\usepackage{amsmath,amssymb,amsthm}
\usepackage{mathtools}
\usepackage[most]{tcolorbox}
\usepackage{xcolor}
\usepackage{tikz}
\usetikzlibrary{arrows.meta,positioning,calc,cd}
\usepackage{booktabs,array,tabularx}
\usepackage{fancyvrb}
\usepackage[hidelinks]{hyperref}
\usepackage{xurl}
\usepackage{enumitem}
\setlist{nosep,leftmargin=*}
\setlength{\emergencystretch}{2em}
\setlength{\parskip}{0.35em}
\numberwithin{equation}{section}

\newcommand{\Set}{\mathbf{Set}}
\newcommand{\PSh}{\mathrm{PSh}}
\newcommand{\Sh}{\mathrm{Sh}}
\newcommand{\Kar}{\operatorname{Kar}}
\newcommand{\Hom}{\operatorname{Hom}}
\newcommand{\End}{\operatorname{End}}
\newcommand{\Aut}{\operatorname{Aut}}
\newcommand{\id}{\mathrm{id}}
\newcommand{\B}{\mathbf{B}}
\newcommand{\y}{\mathbf{y}}
\newcommand{\op}{\mathrm{op}}
\newcommand{\res}[2]{#1|_{#2}}
\newcommand{\q}[1]{«\,#1\,»}
\newcommand{\lean}[1]{{\small\nolinkurl{#1}}}
\newcommand{\lp}{\href{https://github.com/LeanPool/lean-pool/tree/main/LeanPool/ToposTheoreticBridges}{\nolinkurl{ToposTheoreticBridges}}}

% Lean code blocks.
\DefineVerbatimEnvironment{leancode}{Verbatim}{fontsize=\small,frame=leftline,framerule=0.6pt,rulecolor=\color{blue!40},xleftmargin=6pt}

\theoremstyle{plain}
\newtheorem{theorem}{\thmname}[section]
\newtheorem{corollary}[theorem]{\corname}
\newtheorem{proposition}[theorem]{\propname}
\newtheorem{lemma}[theorem]{\lemname}
\newtheorem{bridge}{\bridgename}
\renewcommand{\thebridge}{\Alph{bridge}}
\theoremstyle{definition}
\newtheorem{definition}[theorem]{\defname}
\newtheorem{example}[theorem]{\exname}
\newtheorem{exercise}{\exoname}[section]
\theoremstyle{remark}
\newtheorem*{remark}{\remname}
\tcolorboxenvironment{theorem}{enhanced, breakable, colback=green!2, colframe=green!35!black, boxrule=0.6pt, arc=0pt, left=7pt, right=7pt, top=4pt, bottom=4pt, before skip=12pt, after skip=8pt}
\tcolorboxenvironment{bridge}{enhanced, breakable, colback=blue!3, colframe=blue!35!black, boxrule=0.5pt, arc=0pt, left=7pt, right=7pt, top=4pt, bottom=4pt, before skip=12pt, after skip=8pt}
\tcolorboxenvironment{definition}{enhanced, breakable, colback=black!3, colframe=black!25, boxrule=0.4pt, arc=0pt, left=7pt, right=7pt, top=4pt, bottom=4pt, before skip=10pt, after skip=6pt}
\newcommand{\entry}[1]{\par\smallskip\noindent\textbf{#1}\enspace}

\newcommand{\inwords}{\entry{En mots.}}
\newcommand{\tracename}{Trace du pipeline}
\newcommand{\gen}{\entry{Comment le générateur le trouve.}}
\newcommand{\direct}{\entry{Preuve, sans topos.}}
\newcommand{\status}{\entry{Est-ce nouveau?}}
\newcommand{\leanentry}{\entry{Lean.}}

\title{Les ponts topos-théoriques comme générateur de théorèmes\\[0.3em] \large Nouveaux théorèmes d'un générateur exécuté à la main, décrit comme un pipeline logiciel}
\author{Michel Hua\\ \small Commutator, Paris\\ \small\texttt{michel@commutator.io} · ORCID 0009-0002-8304-0683}
\date{Brouillon, 23 novembre 2026}

\begin{document}
\maketitle
\begin{abstract}
Les \q{ponts} d'O.~Caramello transforment une propriété d'un topos, lue sur deux présentations différentes de ce topos, en un théorème qui relie ces deux présentations. Dans un exposé aux Mines Paris (2024), elle a proposé de confier à une machine la combinaison d'une base de données de telles présentations avec une base de données de telles lectures, afin d'engendrer des théorèmes. Nous exécutons à la main une petite version de ce générateur et la décrivons comme une pipeline logicielle: deux registres en entrée (cinq ponts démontrés en Lean~4 dans la bibliothèque \lp{} de lean-pool; de petites catégories dont les préfaisceaux sont des structures familières), puis jointure, spécialisation, dépliage, réécriture, vérification (en Lean et par une preuve directe) et filtrage par la littérature. Nous ne rapportons que les résultats que nous n'avons pas trouvés publiés: (1) les représentations d'un carquois dans les ensembles vérifient le tiers exclu faible $\neg p\vee\neg\neg p$ exactement quand chaque sommet a au plus une flèche sortante (les automates sur deux lettres ne le vérifient donc pas); (2) deux carquois n'ont des catégories de représentations ensemblistes équivalentes que s'ils sont isomorphes, et les points essentiels sont les sommets; (3) les représentations du carquois cyclique à $n$ sommets ont exactement $2^n+1$ sous-topos, obtenus en inversant un ensemble quelconque de flèches ou en écrasant tout; (4) les points essentiels et les sous-topos des actions de $T_n$, $PT_n$, $I_n$, $O_n$, $M_n(\mathbb F_q)$ sont respectivement au nombre de $n$, $n+1$, $n+1$, $n$, $n+1$ et de $n+1$, $n+2$, $n+2$, $n+1$, $n+2$; (5) un transfert de la condition d'Ore aux coins $eMe$; (6) le tiers exclu faible pour une douzaine de topos combinatoires. Les résultats connus sont consignés à part, pour l'étalonnage, et une dernière section discute les limites de la démarche. Une version anglaise accompagne ce texte.
\end{abstract}

\tableofcontents
\newpage
\section{Le générateur comme pipeline logiciel}\label{sec:intro}

\subsection{Les ponts}

Les \emph{ponts} d'O.~Caramello \cite{caramello-unification,caramello-tst} reposent sur une observation. Un \emph{topos} (une catégorie qui se comporte comme la catégorie des ensembles) peut se présenter de bien des manières qui semblent sans rapport: comme les actions d'un monoïde, comme les représentations d'un carquois, comme les faisceaux sur un espace, comme les modèles d'une théorie logique. Une propriété du topos, par exemple \q{le tiers exclu est vrai}, peut se caractériser sur chaque présentation. Deux caractérisations d'une même propriété, sur deux présentations d'un même topos, donnent un théorème qui relie les deux présentations et où le topos n'apparaît plus. En Lean, deux lemmes \texttt{P E ↔ Q₁ C} et \texttt{P E ↔ Q₂ D}, pour des équivalences \texttt{E ≌ PSh C} et \texttt{E ≌ PSh D}, se composent en \texttt{Q₁ C ↔ Q₂ D}.

Dans un exposé aux Mines Paris, le 17 septembre 2024 \cite{caramello-mines}, Caramello a proposé d'automatiser ce procédé: un assistant de preuve prendrait \q{une base de données d'équivalences de topos} (les \emph{tabliers} des ponts) et \q{une base de données de caractérisations d'invariants topos-théoriques en termes de différentes présentations des topos} (leurs \emph{arches}), et les combinerait \q{pour engendrer des correspondances concrètes entre des «~dépliages~» d'un invariant donné en termes de différentes présentations d'un même topos} \cite[transparent 11]{caramello-mines}. Les résultats, prévient-elle, \q{peuvent être assez «~bizarres~» selon les critères mathématiques habituels}, mais \q{avec un choix soigneux d'équivalences de topos et d'invariants, on peut facilement obtenir des résultats mathématiques intéressants et naturels} \cite[transparent 8]{caramello-mines}. Les citations sont traduites de l'anglais par nos soins.

\subsection{Le pipeline}

Nous faisons tourner un tel générateur à la main, et nous le décrivons comme on décrit un pipeline de données: ses entrées, le rôle de chaque étape, ses sorties, la manière d'en tester les résultats.

\begin{figure}[ht]
\centering
\begin{tikzpicture}[st/.style={draw,rounded corners,align=center,font=\small,minimum height=9mm,text width=25mm},db/.style={draw,align=center,font=\small,minimum height=9mm,text width=25mm,fill=blue!6},ok/.style={draw,rounded corners,align=center,font=\small,minimum height=9mm,text width=25mm,fill=green!8},>=Stealth,node distance=5mm]
\node[db] (decks) at (0,1) {registre des tabliers\\ \footnotesize formes $C$};
\node[db] (arches) at (0,-1) {registre des arches\\ \footnotesize lemmes Lean};
\node[st] (join) at (3.4,0) {1.~jointure\\ \footnotesize tablier $\times$ arche};
\node[st] (spec) at (6.8,0) {2.~spécialiser\\ \footnotesize calculer sur $C$};
\node[st] (unf) at (10.2,0) {3.~déplier\\ \footnotesize termes élémentaires};
\node[st] (rw) at (10.2,-2.2) {4.~réécrire\\ \footnotesize énoncé lisible};
\node[ok] (ver) at (6.8,-2.2) {5.~vérifier\\ \footnotesize Lean + preuve directe};
\node[ok] (nov) at (3.4,-2.2) {6.~filtre de nouveauté\\ \footnotesize littérature};
\node[st] (out) at (0,-3.3) {sortie\\ \footnotesize théorèmes nouveaux};
\node[st] (cal) at (3.4,-4.2) {étalonnage\\ \footnotesize résultats connus};
\draw[->] (decks) -- (join); \draw[->] (arches) -- (join);
\draw[->] (join) -- (spec); \draw[->] (spec) -- (unf); \draw[->] (unf) -- (rw);
\draw[->] (rw) -- (ver); \draw[->] (ver) -- (nov); \draw[->] (nov) -- (out); \draw[->] (nov) -- (cal);
\end{tikzpicture}
\caption{Le générateur, étape par étape. En bleu: les entrées. En vert: les deux étapes de test.}\label{fig:pipeline}
\end{figure}

\begin{description}
\item[Registre des arches.] Cinq lemmes de la forme \texttt{(propriété de PSh C) ↔ (condition sur C)}, démontrés en Lean~4 dans la bibliothèque \lp{} de lean-pool (\S\ref{sec:registry}). C'est sur eux que repose la confiance: tout ce qui suit en est une spécialisation.
\item[Registre des tabliers.] De petites catégories $C$ dont les préfaisceaux forment une structure familière: $\B M$ pour un monoïde $M$ (ses actions), la catégorie libre d'un carquois (ses représentations), $\Delta$ (les ensembles simpliciaux), un préordre (les foncteurs sur ce préordre, ou les faisceaux sur son espace d'Alexandrov). Le livre de La Palme Reyes, Reyes et Zolfaghari \cite{reyes-generic} est un catalogue de tels tabliers, où chaque structure familière est présentée par ses \q{figures génériques}.
\item[1. Jointure.] Associer chaque tablier à chaque arche.
\item[2. Spécialiser.] Appliquer l'arche au tablier et calculer la condition sur $C$: par exemple, la condition d'Ore sur $\B M$ devient $aM\cap bM\neq\emptyset$.
\item[3. Déplier.] Traduire la propriété du topos dans le langage des structures, avec une règle par propriété (le lemme~\ref{lem:neg} pour la négation).
\item[4. Réécrire.] Reformuler la condition sur $C$ dans le langage usuel du domaine (idéaux à droite minimaux, flèches sortantes, etc.). C'est l'étape où il faut un humain, ou un modèle de langage.
\item[5. Vérifier.] Deux tests indépendants: une preuve Lean de l'énoncé spécialisé, comme corollaire de l'arche; et une preuve directe, sans topos, lisible par un étudiant de troisième année de licence.
\item[6. Filtre de nouveauté.] Chercher l'énoncé dans le jeu de données d'appariements de TheoremGraph et par la recherche sémantique de TheoremSearch, puis dans la littérature. Les résultats dont on trouve une publication vont au journal d'étalonnage (\S\ref{sec:calibration}); les autres forment la sortie (\S\ref{sec:machines}--\ref{sec:shapes}).
\end{description}

Chaque théorème ci-dessous est accompagné de sa \emph{trace}, c'est-à-dire de ce que donne chacune des étapes 1 à 6.

\subsection{Résultats}

\begin{center}\small
\begin{tabularx}{\linewidth}{>{\raggedright\arraybackslash}p{4.2cm}>{\raggedright\arraybackslash}X l}
\toprule
tablier & énoncé engendré & \S\\
\midrule
représentations d'un carquois $R$ & tiers exclu faible $\iff$ au plus une flèche sortante par sommet & \ref{sec:machines}\\
automates sur $\Sigma$ & tiers exclu faible $\iff|\Sigma|\le 1$ & \ref{sec:machines}\\
représentations de carquois & catégories équivalentes $\iff$ carquois isomorphes; points essentiels $=$ sommets & \ref{sec:machines}\\
représentations du $n$-cycle & exactement $2^n+1$ sous-topos: inverser un ensemble quelconque de flèches, ou tout écraser & \ref{sec:machines}\\
$T_n$-, $PT_n$-, $I_n$-, $O_n$-, $M_n(\mathbb F_q)$-ensembles & $n$, $n{+}1$, $n{+}1$, $n$, $n{+}1$ points essentiels; $n{+}1$, $n{+}2$, $n{+}2$, $n{+}1$, $n{+}2$ sous-topos & \ref{sec:finite}\\
coins $eMe$, $1\in MeM$ & la condition d'Ore se transfère (aussitôt, par des théorèmes de Morita publiés) & \ref{sec:finite}\\
$M_n(k)$-ensembles; $[\mathrm{Sub}(V),\Set]$ & $M_n(k)$-ensembles $\simeq M_m(k)$-ensembles si et seulement si $n = m$; le treillis des sous-espaces détermine $\dim V$, et le corps lorsque $\dim V\ge 3$ & \ref{sec:linear}\\
$(\mathbb Z/n,\cdot)$-ensembles & $2^{\omega(n)}$ points essentiels, $D(\omega(n))$ sous-topos (nombres de Dedekind) & \ref{sec:linear}\\
faisceaux sur un espace compact séparé $X$ & De Morgan $\iff$ $C(X,\mathbb R)$ complet au sens de Dedekind (par combinaison de théorèmes publiés); pour $X$ métrisable, $\iff$ $X$ fini & \ref{sec:analysis}\\
graphes, ensembles simpliciaux, arbres, etc. & tiers exclu faible, cas par cas (aussitôt, par le pont) & \ref{sec:shapes}\\
\bottomrule
\end{tabularx}
\end{center}

Aucun de ces résultats n'est profond; nous examinons au \S\ref{sec:limits} ce qu'il faut en conclure sur la méthode. Nous écrivons pour un ingénieur qui pratique Lean et qui pense en termes de pipelines: la théorie tient en une courte section, nous citons les énoncés Lean, et nous consacrons la place aux résultats et au travail que demande chaque étape.

\subsection{Plan}

La section~\ref{sec:toolbox} donne les quatre définitions dont le texte a besoin, et le registre: les cinq arches tirées de la littérature, dont viennent la plupart des résultats, et les tabliers. La section~\ref{sec:method} fait du pipeline une recette pratique, avec une table des invariants, une liste de contrôle pour reconnaître les tabliers, un modèle de code Lean et un protocole de recherche bibliographique. Les sections~\ref{sec:machines} à~\ref{sec:shapes} donnent les résultats, tablier par tablier: carquois, monoïdes, analyse, algèbre linéaire et arithmétique, topos combinatoires. La section~\ref{sec:equivalences} fait tourner la même machinerie à rebours pour montrer que des paires de théorèmes classiques de troisième année de licence (Birkhoff et les espaces $T_0$ finis, Cayley et la reconstruction d'un groupe, la théorie de Galois et les revêtements) sont deux lectures d'un même pont. La section~\ref{sec:gluing} ajoute une sixième arche, tirée des notes de Grothendieck sur l'algèbre universelle, et la fait tourner sur les sous-topos essentiels. La section~\ref{sec:calibration} recense les résultats déjà publiés, et la section~\ref{sec:limits} examine les limites de la démarche et les leçons à en tirer pour un générateur automatique. L'annexe~\ref{app:novelty} détaille le filtre de nouveauté et les sources consultées; l'annexe~\ref{sec:lean-lit} donne les fichiers et les déclarations Lean.

Une note associée \cite{bridges-notes} applique le générateur aux notes de Grothendieck: trois arches de plus (deux valeurs, connexe, topos de préfaisceaux d'un préordre), démontrées en Lean et lues sur des tabliers qui reviennent dans ses notes; les précédents qu'on y trouve; enfin le topos des pseudo-corps, dont il donne trois présentations.

\section{Boîte à outils: quatre définitions et les arches}\label{sec:toolbox}

Cette section contient toute la théorie dont la suite a besoin: quatre définitions, et les cinq lemmes qui forment le registre des arches.

\subsection{Formes et préfaisceaux}

Une petite catégorie $C$ est une \emph{forme}. Un \emph{préfaisceau} sur $C$ est un foncteur $C^{\op}\to\Set$, en Lean \texttt{Cᵒᵖ ⥤ Type u}: un ensemble $F(U)$ pour chaque objet et, pour chaque flèche $f : V\to U$, une application de \emph{restriction} $F(U)\to F(V)$, $x\mapsto\res xf$, avec $\res x{\id} = x$ et $\res{(\res xg)}f = \res x{g\circ f}$. Les préfaisceaux sur $C$ forment le topos $\PSh(C)$. Voici les exemples que nous utilisons:

\begin{center}\small
\begin{tabularx}{\linewidth}{>{\raggedright\arraybackslash}p{4.3cm}>{\raggedright\arraybackslash}X}
\toprule
forme $C$ & préfaisceaux sur $C$\\
\midrule
$\B M$: un objet, flèches $=$ éléments d'un monoïde $M$, $g\circ f = gf$ & $M$-ensembles à droite: un ensemble $X$ muni de $x\cdot m$, avec $x\cdot 1 = x$, $(x\cdot g)\cdot f = x\cdot gf$; dans Mathlib, \texttt{MulAction Mᵐᵒᵖ X}\\
$\B(\Sigma^*)$, $\Sigma^*$ les mots sur un alphabet & automates déterministes: états $X$, transitions $x\cdot a$ (sans état initial ni états finaux)\\
$\B(\mathbb N,+)$ & systèmes dynamiques discrets: un ensemble muni d'une application $\tau : X\to X$\\
la catégorie libre d'un carquois & représentations du carquois opposé: un ensemble par sommet, une application par flèche\\
$V\rightrightarrows E$ & graphes orientés\\
$\Delta$, $\Delta_+$ & ensembles simpliciaux, ensembles simpliciaux augmentés\\
$\omega = (\mathbb N,\le)$ & suites $X_0\leftarrow X_1\leftarrow\cdots$ (le \q{topos des arbres} de la récursion gardée)\\
\bottomrule
\end{tabularx}
\end{center}

Pour un monoïde, la restriction de $x$ le long de $m$ est $x\cdot m$.

\subsection{Pièges et négation}

Un \emph{sous-préfaisceau} $A\subseteq F$, que nous appelons aussi un \emph{piège}, est une famille de sous-ensembles stable par restriction. Pour un $M$-ensemble, c'est une partie stable par l'action. Les sous-préfaisceaux sont les sous-objets du topos et jouent le rôle de propositions. Leur négation n'est pas le complémentaire ensembliste (qui n'est en général pas un piège), mais le plus grand sous-préfaisceau disjoint de $A$:

\begin{lemma}\label{lem:neg}
On a $x\in\neg A$ si et seulement si aucune restriction de $x$ n'appartient à $A$. Par conséquent, $x\in\neg\neg A$ si et seulement si toute restriction de $x$ admet une restriction ultérieure dans $A$.
\end{lemma}

\begin{proof}
L'ensemble de ces $x$ est stable par restriction, disjoint de $A$, et contient tout piège disjoint de $A$. La seconde phrase est la première appliquée à $\neg A$.
\end{proof}

\begin{definition}
Le topos $\PSh(C)$ est \emph{booléen} si $A\cup\neg A = F$ pour tout $A\subseteq F$ (tiers exclu), et \emph{de De Morgan} si $\neg A\cup\neg\neg A = F$ pour tout $A\subseteq F$ (tiers exclu faible). D'après le lemme~\ref{lem:neg}, De Morgan signifie: \emph{pour tout préfaisceau $F$, tout $A\subseteq F$ et tout élément $x$, ou bien aucune restriction de $x$ n'appartient à $A$, ou bien toute restriction de $x$ admet une restriction ultérieure dans $A$.}
\end{definition}

Informellement, De Morgan est une propriété de \emph{stabilité}: \q{$A$ est accessible} ne peut pas se perdre en restreignant davantage. Être booléen est plus fort: \q{$x\in A$} est décidé.

\subsection{La condition d'Ore}

La catégorie $C$ vérifie la \emph{condition d'Ore} si toute paire de flèches de même but $f : Y\to X$, $g : Z\to X$ se complète en un carré commutatif $f\circ h = g\circ k$. Pour $\B M$: pour tous $a,b\in M$, il existe $c,d$ avec $ac = bd$, c'est-à-dire $aM\cap bM\neq\emptyset$ (deux éléments quelconques ont un multiple à droite commun). Elle est vraie pour les groupes, les monoïdes commutatifs, les monoïdes à zéro, les catégories ayant un objet initial ou des produits fibrés, et les chaînes.

\subsection{Les cinq arches}

La bibliothèque \lp{} démontre en Lean~4 les cinq ponts suivants. Nous citons les déclarations principales; les deux premières sont données avec leur énoncé exact.

\begin{bridge}[Booléen]\label{br:bool}
$\PSh(C)$ est booléen si et seulement si $C$ est un groupoïde.
\begin{leancode}
theorem isBooleanTopos_presheaf_iff_isGroupoid (C : Type u) [SmallCategory C] :
    IsBooleanTopos (Cᵒᵖ ⥤ Type u) ↔ IsGroupoid C
\end{leancode}
\end{bridge}

\begin{bridge}[De Morgan]\label{br:dm}
$\PSh(C)$ est de De Morgan si et seulement si $C$ vérifie la condition d'Ore.
\begin{leancode}
theorem isDeMorganTopos_presheaf_iff_rightOreCondition (C : Type u)
    [SmallCategory C] :
    IsDeMorganTopos (Cᵒᵖ ⥤ Type u) ↔ GrothendieckTopology.RightOreCondition C
\end{leancode}
Sa preuve est courte. Si $\res xf\in A$ et si $g$ est une autre flèche de but $X$, on complète le carré $f\circ h = g\circ k$; alors $\res{(\res xg)}k = \res{(\res xf)}h\in A$. Réciproquement, on teste sur le préfaisceau $\y X = \Hom(-,X)$, le piège engendré par $f$, et l'élément $\id_X$.
\end{bridge}

\begin{bridge}[Morita]\label{br:morita}
On a $\PSh(C)\simeq\PSh(D)$ si et seulement si les \emph{enveloppes de Karoubi} $\Kar(C)$ et $\Kar(D)$ sont équivalentes. $\Kar(C)$ ajoute un objet pour chaque flèche idempotente $e = e\circ e$; pour $\B M$, ses objets sont les idempotents de $M$, et $e\cong f$ dans $\Kar(\B M)$ si et seulement s'il existe $a\in fMe$, $b\in eMf$ avec $ba = e$ et $ab = f$. Lean: \lean{moritaEquivalentCat_iff_karoubi}; corollaires \lean{isGroupoid_iff_of_moritaEquivalentCat}, \lean{rightOreCondition_iff_of_moritaEquivalentCat}.
\end{bridge}

\begin{bridge}[Points essentiels]\label{br:pts}
Les foncteurs $\PSh(C)\to\Set$ qui préservent toutes les limites et toutes les colimites (les \emph{points essentiels}, ou \q{foncteurs fibres}) forment une catégorie équivalente à $\Kar(C)^{\op}$. Pour $\B M$, ce sont les foncteurs $X\mapsto Xe = \{x : x\cdot e = x\}$, pour $e$ idempotent, deux d'entre eux étant isomorphes si et seulement si $e\cong f$ dans $\Kar(\B M)$. Lean: \lean{essentialPointEquivKaroubiOp}, \lean{isEssentialPoint_iff}.
\end{bridge}

\begin{bridge}[Sous-topos]\label{br:sub}
Les sous-topos de $\PSh(C)$ (sous-catégories pleines stables par limites dont l'inclusion admet un adjoint à gauche qui préserve les limites finies) correspondent, en renversant l'ordre, aux \emph{topologies de Grothendieck} sur $C$: des choix de cribles \q{couvrants} sur chaque objet, stables par changement de base et transitifs. Lean: \lean{subtoposOrderIso}, \lean{exists_orderIso_ici_of_sheafEquiv}.
\end{bridge}

C'est tout. Chaque énoncé qui suit est l'un de ces cinq lemmes, spécialisé à une forme et déplié à l'aide du lemme~\ref{lem:neg}.

\subsection{Trois arches de plus}\label{sec:morearches}

Trois autres arches ont leur place dans le registre. Elles n'ont été jointes pour l'instant qu'aux tabliers des notes de Grothendieck, ci-dessous; aucun résultat des sections suivantes n'en provient.

\begin{bridge}[Deux valeurs]\label{br:twoval}
Les sous-objets de l'objet final de $\PSh(C)$ (ses \q{ouverts}) correspondent aux cribles de $C$ \cite[IV~8.4.4]{sga4}. Ainsi $\PSh(C)$ n'a pas d'autre sous-objet de l'objet final que $0$ et $1$ si et seulement si, pour deux objets quelconques $a$, $b$ de $C$, il existe une flèche $a\to b$. Lean: \lean{forall_subfunctor_terminal_eq_bot_or_eq_top_iff}.
\end{bridge}

\begin{bridge}[Connexe]\label{br:conn}
$\PSh(C)$ est connexe non vide au sens de \cite[IV~8.7]{sga4} (toute partition de l'objet final a exactement une partie non vide) si et seulement si $C$ est connexe. En termes de sous-objets de l'objet final: $C$ est non vide et les seuls sous-objets complémentés de l'objet final sont $0$ et $1$. Lean: \lean{isConnected_iff_forall_isCompl}.
\end{bridge}

\begin{bridge}[Spatial]\label{br:spatial}
$\PSh(C)$ est équivalent au topos de préfaisceaux d'un préordre, donc au topos des faisceaux sur son espace d'Alexandrov, si et seulement si $C$ est mince. Grothendieck l'énonce comme exercice \cite[p.~26]{grothendieck-161-3}; voir le \S\ref{sec:equivalences}. Lean: \lean{exists_preorder_moritaEquivalentCat_iff_isThin}.
\end{bridge}

\subsection{Les tabliers des notes de Grothendieck}\label{sec:grdecks}

Trois tabliers reviennent dans les notes de travail de Grothendieck, et les trois arches ci-dessus les lisent ainsi (Lean: \lean{GrothendieckDecks.lean}).
\begin{itemize}
\item La catégorie simpliciale $\Delta$, la première catégorie test de \emph{Pursuing Stacks} (dossiers 108 et 118 du fonds): les ensembles simpliciaux forment un topos à deux valeurs et connexe, puisqu'il existe une application constante entre deux objets quelconques de $\Delta$, et ce n'est pas le topos de préfaisceaux d'un préordre, puisque les deux sommets de $[1]$ sont deux flèches $[0]\to[1]$ (\lean{simplexCategory_twoValued}, \lean{simplexCategory_isConnected}, \lean{simplexCategory_not_spatial}).
\item Un groupe $G$ vu comme groupoïde à un objet, dont le topos de préfaisceaux est $\B G$ (dossier 161-6, p.~22, où Grothendieck ramène une question à ce cas): à deux valeurs, comme le note \cite[IV~8.4.3]{sga4} pour tout monoïde, connexe, et topos de préfaisceaux d'un préordre si et seulement si $G$ est trivial (\lean{singleObj_twoValued}, \lean{singleObj_isConnected}, \lean{singleObj_spatial_iff}).
\item Un ensemble ordonné $I$, dont le topos de préfaisceaux $\widehat I$ est $\mathrm{Top}(I)$ (dossier 37, p.~48; dossier 161-3, p.~26): à deux valeurs si et seulement si deux éléments quelconques sont comparables dans les deux sens, c'est-à-dire si et seulement si son quotient ordonné a au plus un point (\lean{preorder_twoValued_iff}).
\end{itemize}
Ce sont des corollaires immédiats, consignés parce qu'ils sont les premières entrées du registre tirées du fonds.

\section{Une méthode pratique pour produire des théorèmes}\label{sec:method}

Cette section fait de la pipeline de la figure~\ref{fig:pipeline} une recette qu'un étudiant, ou un programme, peut suivre. Chaque étape a une entrée, une sortie et un contrôle.

\subsection{Étape 1: choisir un invariant lisible des deux côtés}

On part du registre des arches. Pour les topos de préfaisceaux, chaque arche vient avec une lecture sur la forme $C$ et une \emph{règle de dépliage} qui traduit la propriété du topos en une phrase sur les structures:

\begin{center}\small
\begin{tabularx}{\linewidth}{>{\raggedright\arraybackslash}p{2.6cm}>{\raggedright\arraybackslash}p{3.2cm}>{\raggedright\arraybackslash}X>{\raggedright\arraybackslash}X}
\toprule
invariant & lecture sur $C$ & sur les $M$-ensembles à droite & sur les représentations d'un carquois\\
\midrule
booléen & groupoïde & $M$ est un groupe & aucune flèche\\
De Morgan & condition d'Ore & $aM\cap bM\neq\emptyset$ & $\le 1$ flèche sortante par sommet\\
classe de Morita & $\Kar(C)$ & coins $eMe$, $1\in MeM$ & le carquois lui-même\\
points essentiels & objets de $\Kar(C)$ & idempotents modulo $\mathcal D$ & sommets\\
sous-topos & topologies sur $C$ & idéaux idempotents ($M$ fini) & par ex.\ $2^n+1$ pour le $n$-cycle\\
\bottomrule
\end{tabularx}
\end{center}

\emph{Contrôle:} l'invariant doit être invariant par équivalence de catégories. Les propriétés énoncées à l'aide de sous-objets, de limites et de colimites le sont; celles qui font référence à une présentation choisie (\q{la forme a trois objets}) ne le sont pas.

\subsection{Étape 2: reconnaître un tablier}

Une structure est un topos de préfaisceaux quand elle a l'anatomie suivante:
\begin{enumerate}
\item un ensemble de données en chaque \q{lieu} (un objet de $C$);
\item une manière de transporter les données le long de \q{relations} entre les lieux (les applications de restriction), contravariante et compatible à la composition;
\item aucune équation ni condition de recollement supplémentaire.
\end{enumerate}
Si (3) est en défaut parce que les données doivent se recoller (fonctions sur des ouverts, sections d'un revêtement), la structure est un topos de faisceaux sur un site; l'arche du pont~\ref{br:sub} s'applique encore, les quatre autres demandent des versions pour les sites. Tabliers typiques: actions d'un monoïde ($C = \B M$), représentations d'un carquois ($C$ = catégorie libre), objets simpliciaux ($C = \Delta$), familles indexées par un préordre ($C = P^{\op}$), modèles d'une théorie de Horn ($C$ = opposée des modèles de présentation finie). \emph{Contrôle:} écrire $C$ explicitement et vérifier, sur un ou deux exemples, qu'un préfaisceau sur $C$ est bien la structure visée; le côté de l'action et le sens des flèches sont les sources d'erreur habituelles.

\subsection{Étape 3: spécialiser l'arche}

On calcule la lecture sur $C$. Deux heuristiques couvrent la plupart des cas:
\begin{itemize}
\item \emph{Tester sur les représentables.} Une propriété de tous les préfaisceaux se décide en général sur les préfaisceaux représentables $\y U = \Hom(-,U)$ et leurs cribles; c'est ainsi que sont démontrés les ponts~\ref{br:bool} et~\ref{br:dm}, et ainsi que l'on construit des contre-exemples (figure~\ref{fig:branch}).
\item \emph{Utiliser une forme normale pour les flèches.} Dans une catégorie libre, les flèches sont des chemins; dans $\B M$, des éléments; dans un préordre, il y en a au plus une. La condition devient un énoncé sur des mots, des éléments ou des inégalités.
\end{itemize}
\emph{Contrôle:} énumérer les petits cas par machine quand c'est possible (petits monoïdes, petits carquois) et comparer avec la condition spécialisée.

\subsection{Étape 4: déplier et réécrire}

On traduit la propriété du topos avec la règle de dépliage (le lemme~\ref{lem:neg} pour la négation), puis on traduit la condition sur $C$ dans le vocabulaire usuel du domaine (idéaux à droite minimaux, degré sortant, rang). On garde les deux formes: la forme brute est celle que Lean démontre, la forme réécrite est celle que comprend un lecteur.

\subsection{Étape 5: vérifier deux fois}

\emph{En Lean.} Un énoncé engendré est un corollaire en deux lignes de l'arche, une fois la condition sur $C$ démontrée. Le schéma, tiré de la bibliothèque:
\begin{leancode}
theorem isDeMorganTopos_singleObj_iff :
    IsDeMorganTopos ((SingleObj M)ᵒᵖ ⥤ Type) ↔ ∀ a b : M, ∃ c d : M, a * c = b * d := by
  rw [isDeMorganTopos_presheaf_iff_rightOreCondition, rightOreCondition_singleObj_iff]

theorem not_isDeMorganTopos_freeMonoid_two :
    ¬ IsDeMorganTopos ((SingleObj (FreeMonoid (Fin 2)))ᵒᵖ ⥤ Type) := by
  rw [isDeMorganTopos_singleObj_iff]
  -- reste : x et y n'ont pas de multiple à droite commun dans le monoïde libre
\end{leancode}
Tout le travail est dans le lemme sur $C$ (ici \lean{rightOreCondition_singleObj_iff}), qui relève de l'algèbre ou de la combinatoire ordinaires. \emph{À la main.} On écrit une preuve directe qui n'utilise pas les topos; si elle est bien plus difficile que le pont, le pont a apporté quelque chose.

\subsection{Étape 6: filtrer la nouveauté}

On commence par un passage mécanique. On résume en une phrase l'énoncé Lean du théorème réécrit et l'on cherche ce résumé dans le jeu de données d'appariements de TheoremGraph (fiche du jeu de données, Hugging Face, consultée le 23 novembre 2026) \cite{theoremgraph-matching}. Chacune de ses 100\,831 lignes associe une déclaration Lean, de Mathlib ou de l'un de 25 autres projets, à l'énoncé d'arXiv le plus proche, trouvé par similarité cosinus (au moins 0,80) entre plongements de résumés, et classé \emph{exact}, \emph{inexact} ou \emph{faux} par un modèle de langage; son juge principal tient 47\,952 paires pour des appariements. Une réponse positive établit qu'un jumeau formel existe et désigne l'article qui en donne la version informelle; le résultat va alors au journal d'étalonnage, avec une référence précise. Une réponse négative ne prouve rien: un énoncé informel sans déclaration Lean ne figure pas dans le jeu, les étiquettes sont celles d'un modèle et ne sont pas vérifiées, et le texte informel n'est fourni que pour les paires dont la source arXiv autorise la redistribution. La recherche à la main reste donc nécessaire. Nous n'avons fait ce passage que sur huit résultats (les énoncés sur les carquois, le carquois cyclique, Morita, les coins, $T_n$, $M_n(k)$, $(\mathbb Z/n,\cdot)$ et les compacts séparés; ni les automates, ni les topos combinatoires, ni les treillis de sous-espaces), par mots-clés, dans le service de recherche public du jeu, et aucune des paires renvoyées ne concernait l'un d'eux. On s'y attendait, et cela ne dit rien de la nouveauté: le côté Lean du jeu ne contient aucune déclaration de Lean Pool, ni aucune sur les topos de De Morgan, et une recherche par mots-clés classe selon les mots, non selon la similarité des résumés. Pour le côté informel, nous avons ensuite soumis les dix lignes du tableau des résultats à la recherche sémantique de TheoremSearch \cite{theoremsearch}, qui plonge une description en langue naturelle et classe quelque 8 millions d'énoncés tirés d'arXiv et d'autres sources selon la similarité de leurs résumés (API publique, 23 novembre 2026, six résultats par requête). Aucun des énoncés renvoyés ne formule l'un des dix. Trois requêtes en ont renvoyé les ingrédients, que l'article signale déjà: le théorème fondamental de la géométrie projective, pour les treillis de sous-espaces; pour le résultat sur les compacts séparés, la caractérisation des topos localiques de De Morgan (celle de Johnstone, sous la forme que lui donne Caramello) et celle des compacts séparés dont $C(X,\mathbb R)$ est complet au sens de Dedekind comme les espaces stoniens, en deux énoncés distincts; pour les coins, des énoncés sur les anneaux où une propriété de $R$ passe à $eRe$ quand $ReR = R$, mais pas l'énoncé sur les monoïdes. Les autres requêtes ont renvoyé des énoncés sur les carquois, les cycles, les monoïdes ou les diviseurs, sans rapport avec nos résultats. Avec six résultats par requête, la recherche reste peu profonde, et un énoncé formulé dans un autre vocabulaire peut lui échapper; ce passage resserre donc la recherche à la main sans la remplacer. Les appels, leurs paramètres et les réponses du 23 novembre 2026 sont dans \texttt{docs/bridges/novelty/} (\texttt{theoremsearch.py} et \texttt{theoremsearch-2026-11-23.json}): une requête \texttt{POST https://api.theoremsearch.com/search} par résultat, avec la description comme \texttt{query}, \texttt{n\_results} fixé à 6, sans filtre ni pondération par les citations.

On cherche ensuite, dans cet ordre: l'\emph{Elephant} de Johnstone; les articles de l'auteur de l'arche (ici Caramello); la littérature récente sur le tablier (ici Hemelaer, Rogers, Pirashvili pour les actions de monoïdes); la littérature classique du domaine de l'énoncé réécrit (demi-groupes, carquois, topologie); les serveurs de prépublications et le nLab. Il faut chercher l'énoncé \emph{réécrit}, dans le vocabulaire de son domaine, et pas seulement l'énoncé topos-théorique: trois de nos candidats (le critère pour les monoïdes finis, les sous-topos des $\mathbb N$-ensembles, les points essentiels comme classes d'idempotents) étaient publiés dans le langage des demi-groupes ou des actions de monoïdes, alors que nous les avions crus nouveaux. Les énoncés trouvés sont consignés au journal d'étalonnage (\S\ref{sec:calibration}).

\subsection{Étape 7: décider si c'est intéressant}

Un énoncé engendré mérite d'être gardé quand (a) il relie deux vocabulaires que l'on ne relie pas d'habitude, (b) sa forme réécrite est courte, (c) sa preuve directe n'est pas immédiate \emph{ou} la question à laquelle il répond n'avait pas été posée. Un candidat qui échoue aux trois critères est écarté; le nombre de candidats écartés n'a pas été consigné (\S\ref{sec:calibration}), et le classement des candidats est une limite de la méthode (\S\ref{sec:limits}).

\section{Carquois}\label{sec:machines}

Un \emph{carquois} $R$ est un graphe orienté (flèches multiples et boucles permises). Une \emph{représentation} de $R$ dans les ensembles donne un ensemble $F(v)$ pour chaque sommet et une application $F(a) : F(v)\to F(w)$ pour chaque flèche $a : v\to w$, sans équations; les morphismes sont les familles d'applications qui commutent aux $F(a)$. On note $\mathrm{Rep}(R)$ cette catégorie. C'est le topos des préfaisceaux sur la catégorie libre $C_R = \mathrm{Paths}(R^{\op})$, dont les flèches sont les chemins du carquois opposé: un chemin est la seule forme de composé possible, et l'équation $f\circ h = g\circ k$ dans $C_R$ est une égalité de chemins.

Ce seul tablier, passé par les cinq arches, donne quatre énoncés. Deux faits sur les catégories libres servent partout: la longueur des chemins est additive, donc un chemin idempotent ($p\circ p = p$) est de longueur $0$ et c'est une identité; un chemin inversible est lui aussi de longueur $0$.

\subsection{De Morgan: le déterminisme}

On dit qu'un élément $x\in F(v)$ est \emph{évitant} pour une sous-représentation $A$ si aucune image de $x$ le long d'un chemin n'appartient à $A$, et \emph{attiré} par $A$ si toute image de $x$ admet une image ultérieure dans $A$.

\begin{theorem}\label{th:quivers}
Le topos $\mathrm{Rep}(R)$ est de De Morgan (pour toute sous-représentation, tout élément est évitant ou attiré) si et seulement si chaque sommet de $R$ a au plus une flèche sortante. C'est un topos booléen si et seulement si $R$ n'a aucune flèche.
\end{theorem}

\begin{trace}
\item[tablier] $\mathrm{Rep}(R) = \PSh(\mathrm{Paths}(R^{\op}))$.
\item[arche] Pont~\ref{br:dm} (De Morgan $\iff$ Ore) et pont~\ref{br:bool} (booléen $\iff$ groupoïde).
\item[spécialiser] $f\circ h = g\circ k$ est une égalité de chemins, donc l'un de $f$, $g$ est un segment final de l'autre; Ore sur $\mathrm{Paths}(R^{\op})$ $\iff$ de deux chemins aboutissant à un même sommet de $R^{\op}$, l'un est un segment final de l'autre $\iff$ aucun sommet de $R^{\op}$ n'a deux flèches entrantes. Une catégorie libre est un groupoïde si et seulement si elle n'a aucune flèche.
\item[déplier] Lemme~\ref{lem:neg}: $\neg A$ = évitants, $\neg\neg A$ = attirés.
\item[réécrire] deux flèches entrantes dans $R^{\op}$ = deux flèches sortantes dans $R$.
\item[vérifier] preuve directe ci-dessous; Lean: \lean{isDeMorganTopos_paths_iff} et \lean{isBooleanTopos_paths_iff}, à partir de \lean{rightOreCondition_paths_iff} (dans la convention de Lean, les préfaisceaux sur \texttt{Paths Q} sont les représentations du carquois opposé, et la condition s'écrit \q{au plus une flèche entrante en chaque sommet de $Q$}) (branche \texttt{mycaule/topos-bridges-new}, commit \texttt{463752c}).
\item[nouveauté] un cas particulier est publié (le carquois $\bullet\leftarrow\bullet\to\bullet$, \cite[\S4.1.4]{caramello-chamoun}); la caractérisation n'a pas été trouvée.
\end{trace}

\direct{} Si chaque sommet a au plus une flèche sortante, les chemins issus de $v$ forment une seule suite $p_0\subset p_1\subset\cdots$, chacun prolongeant le précédent; si $F(p_k)(x)\in A$, alors, depuis toute image $F(p_j)(x)$, l'image ultérieure $F(p_{\max(j,k)})(x)$ appartient à $A$. Réciproquement, pour deux flèches $a : v\to w$ et $b : v\to w'$, soit $F(u)$ l'ensemble des chemins de $v$ à $u$ (les applications prolongent les chemins), $A$ les chemins qui commencent par $a$, et $x$ le chemin vide en $v$ (figure~\ref{fig:branch}): il n'est ni évitant ni attiré. Pour le cas booléen, si $a : v\to w$ est une flèche, on prend $F(v) = F(w) = \{\ast\}$, $A(w) = \{\ast\}$, $A(v) = \emptyset$.\qed

\begin{figure}[ht]
\centering
\begin{tikzpicture}[vx/.style={circle,draw,inner sep=1.5pt,minimum size=6mm,font=\small},>=Stealth]
\node[vx] (v) at (0,0) {$v$}; \node[vx,fill=red!12] (w) at (2.6,0.8) {$w$}; \node[vx] (w2) at (2.6,-0.8) {$w'$};
\draw[->] (v) -- node[above,font=\small]{$a$} (w); \draw[->] (v) -- node[below,font=\small]{$b$} (w2);
\node[align=left,font=\small] at (7.6,0) {$F(u)$ = chemins de $v$ à $u$; $A$ = chemins commençant par $a$;\\ le chemin vide en $v$ n'est ni évitant ni attiré.};
\end{tikzpicture}
\caption{Un sommet qui se ramifie met en défaut le tiers exclu faible.}\label{fig:branch}
\end{figure}

\begin{corollary}\label{co:automata}
Les actions à droite du monoïde libre $\Sigma^*$ (automates déterministes sur $\Sigma$, sans état initial ni états finaux) forment un topos de De Morgan si et seulement si $|\Sigma|\le 1$, et un topos booléen si et seulement si $\Sigma = \emptyset$.
\end{corollary}

\begin{proof}
Une action de $\Sigma^*$ est une représentation du carquois à un sommet et $|\Sigma|$ boucles.
\end{proof}

\begin{trace}
\item[vérifier] Lean: \lean{isDeMorganTopos_singleObj_freeMonoid_unit} (une lettre), \lean{not_isDeMorganTopos_freeMonoid_two} (deux lettres), à partir de \lean{isDeMorganTopos_singleObj_iff}.
\item[nouveauté] le topos des automates est classique; cette propriété n'a pas été trouvée dans la littérature.
\end{trace}

\subsection{Morita: le carquois est déterminé par ses représentations}

\begin{theorem}\label{th:quivermorita}
Deux carquois $R$ et $R'$ ont des catégories de représentations ensemblistes équivalentes si et seulement s'ils sont isomorphes.
\end{theorem}

\begin{trace}
\item[tablier] $\mathrm{Rep}(R)\simeq\mathrm{Rep}(R')$, c'est-à-dire $\PSh(C_R)\simeq\PSh(C_{R'})$.
\item[arche] Pont~\ref{br:morita}: $\iff\Kar(C_R)\simeq\Kar(C_{R'})$.
\item[spécialiser] une catégorie libre est karoubienne (les chemins idempotents sont des identités), donc $\Kar(C_R)\simeq C_R$; elle est squelettique et sans isomorphisme autre que les identités, donc une équivalence $C_R\simeq C_{R'}$ est un isomorphisme de catégories.
\item[réécrire] un isomorphisme de catégories libres envoie les flèches indécomposables (chemins de longueur $1$) sur des flèches indécomposables; il est donc induit par un isomorphisme $R^{\op}\cong R'^{\op}$.
\item[vérifier] preuve directe dans la trace; Lean: \lean{moritaEquivalentCat_paths_iff}, \lean{nonempty_representations_equiv_iff}, à partir de \lean{paths_eq_id_of_idempotent}, \lean{paths_nonempty_iso_iff} et \lean{nonempty_paths_equiv_iff} (une équivalence de catégories libres induit un isomorphisme de carquois, via \lean{pathLength_eq_one_iff_indecomposable}) (branche \texttt{mycaule/topos-bridges-quivers}, commit \texttt{2a32779}).
\item[nouveauté] non trouvé comme énoncé sur les carquois; chaque étape est publiée (Morita par l'enveloppe de Karoubi \cite[Lemme~1.1.3]{rogers-thesis}, d'après \cite[A1.1.9--10]{elephant}).
\end{trace}

\inwords{} Pour les représentations d'algèbres \emph{linéaires}, des carquois très différents peuvent avoir des catégories (dérivées) de représentations équivalentes; dans les ensembles, ce phénomène n'existe pas: la catégorie des représentations ensemblistes garde la mémoire du carquois.

\subsection{Points essentiels: les sommets}

\begin{theorem}\label{th:quiverpoints}
Les points essentiels de $\mathrm{Rep}(R)$ (foncteurs vers les ensembles qui préservent toutes les limites et colimites) forment une catégorie isomorphe à la catégorie libre $\mathrm{Paths}(R)$. À isomorphisme près, ce sont les évaluations $F\mapsto F(v)$ aux sommets, deux d'entre elles n'étant isomorphes que si les sommets sont égaux.
\end{theorem}

\begin{trace}
\item[arche] Pont~\ref{br:pts}: points essentiels $\simeq\Kar(C_R)^{\op}\simeq C_R^{\op} = \mathrm{Paths}(R^{\op})^{\op}\cong\mathrm{Paths}(R)$.
\item[vérifier] Lean: \lean{essentialPointPathsEquiv}, \lean{nonempty_evaluationPoint_iso_iff}, \lean{exists_iso_evaluationPoint} (branche \texttt{mycaule/topos-bridges-quivers}, commit \texttt{2a32779}).
\item[nouveauté] le fait général (une petite catégorie se retrouve, à complétion idempotente près, à partir des points essentiels de ses préfaisceaux) est publié \cite[\S4.2.1]{rogers-thesis}; la version pour les carquois n'a pas été trouvée.
\end{trace}

\subsection{Sous-topos du carquois cyclique}

Soit $Z_n$ le carquois cyclique de sommets $\mathbb Z/n$, avec une flèche $a_i : i\to i+1$ pour chaque $i$. Une représentation est un cycle d'applications $F(0)\to F(1)\to\cdots\to F(n-1)\to F(0)$. Pour $n = 1$, c'est un ensemble muni d'un endomorphisme.

\begin{theorem}\label{th:cycle}
Le topos $\mathrm{Rep}(Z_n)$ a exactement $2^n + 1$ sous-topos: pour chaque ensemble $S$ de flèches, la sous-catégorie pleine des représentations dans lesquelles les applications $F(a)$, $a\in S$, sont bijectives; et le sous-topos dégénéré.
\end{theorem}

\begin{trace}
\item[tablier] $\mathrm{Rep}(Z_n) = \PSh(C)$, $C = \mathrm{Paths}(Z_n^{\op})$.
\item[arche] Pont~\ref{br:sub}: sous-topos $\leftrightarrow$ topologies de Grothendieck sur $C$.
\item[spécialiser] classer les topologies sur $C$ (ci-dessous): la topologie dégénérée, ou une par partie de l'ensemble des sommets.
\item[déplier] identifier les faisceaux: une représentation est un faisceau si et seulement si les flèches issues des sommets du complémentaire de la partie sont bijectives.
\item[vérifier] preuve directe ci-dessous; Lean: \lean{card_grothendieckTopology} (exactement $2^n+1$ topologies), \lean{topologyEquiv} (une bijection avec \texttt{Option (Set (ZMod n))}), \lean{eq_top_or_eq_topologyOf} et \lean{CyclicSite.card_grothendieckTopology_one} pour $n = 1$; la description des faisceaux de chaque topologie n'est pas formalisée (branche \texttt{mycaule/topos-bridges-quivers}, commit \texttt{2a32779}).
\item[nouveauté] pour $n = 1$, le résultat découle de la description par Pirashvili des sous-topos des $M$-ensembles pour $M$ commutatif \cite[Cor.~4.3.11]{pirashvili-jpaa}; pour $n\ge 2$, il n'a pas été trouvé. Pour les carquois \emph{sans} cycle orienté, les topologies de Grothendieck sur la catégorie des chemins sont connues: il y en a $2^m$ pour $m$ sommets, une par ensemble de sommets \cite[Cor.~4.3]{schmid} (d'après une note non publiée de Murfet); le cas cyclique est d'une autre nature.
\end{trace}

\direct{} Dans $C$, pour chaque sommet $v$ et chaque $\ell\ge 0$, il y a exactement un chemin de longueur $\ell$ qui aboutit à $v$ (il part de $v+\ell$). Un crible sur $v$ est stable par allongement, c'est donc $\emptyset$ ou $S_{v,k} = \{\text{chemins de longueur}\ge k\}$ pour un $k\ge 0$; $S_{v,0}$ est le crible maximal. Le changement de base de $S_{v,k}$ le long du chemin de longueur $m$ qui aboutit à $v$ donne $S_{v+m,\max(k-m,0)}$.

Soit $J$ une topologie. Si $\emptyset$ couvre un sommet, il couvre tous les sommets (on tire en arrière le long d'un chemin, qui existe car le cycle est fortement connexe), et $J$ est dégénérée. Sinon, les cribles couvrants sont stables par passage à un crible plus grand, donc $J(v) = \{S_{v,k} : k\le K(v)\}$ pour un $K(v)\in\{0,1,2,\dots,\infty\}$. Les axiomes deviennent:
\begin{itemize}
\item \emph{stabilité}: $K(v+1)\ge K(v)-1$;
\item \emph{transitivité}: si $1\le K(v) = k<\infty$, alors $K(v+k) = 0$. (On prend le crible couvrant $S_{v,k}$ et $T = S_{v,k+1}$: les changements de base de $T$ le long des chemins de $S_{v,k}$ sont les $S_{v+r,k+1-r}$, $r\ge k$, tous maximaux sauf $S_{v+k,1}$; si $K(v+k)\ge 1$, ils couvrent tous, donc $T$ couvre, ce qui contredit $K(v) = k$. Réciproquement, quand $K(v+k) = 0$, la transitivité n'impose jamais de nouveau crible couvrant.)
\end{itemize}
Ces deux conditions disent exactement que $K(v)$ est la distance de $v$ au prochain élément de $Z = K^{-1}(0)$ en avançant ($\infty$ si $Z = \emptyset$): une valeur $k\ge 1$ doit être suivie de $k-1$, et une suite $k, k-1,\dots,1$ se termine par un $0$. Les topologies non dégénérées correspondent donc aux parties $Z\subseteq\mathbb Z/n$, et il y en a $2^n$. Enfin, une famille compatible sur $S_{v,1}$ est un élément de $F(v+1)$, et son recollement est un antécédent par $F(a_v)$; donc dire que $S_{v,1}$ couvre (c'est-à-dire $v\notin Z$) revient à demander que $F(a_v)$ soit bijective, et les cribles plus longs n'ajoutent rien. Le sous-topos associé à $Z$ est formé des représentations dans lesquelles les flèches issues des sommets hors de $Z$ sont bijectives; on pose $S = \{a_v : v\notin Z\}$.\qed

\inwords{} Toute manière de \q{localiser} un cycle d'applications consiste à inverser certaines d'entre elles, ou à tout écraser. Pour $n = 1$: les ensembles munis d'un endomorphisme, les ensembles munis d'un automorphisme, et le topos dégénéré. Pour un carquois acyclique, la réponse est différente \cite{schmid}: le topos de Sierpiński $\mathrm{Rep}(\bullet\to\bullet)$ a $2^2 = 4$ sous-topos, dont l'un (\q{$F(\text{but})$ est un point}) n'inverse aucune flèche. Sur un cycle, la connexité forte exclut de tels sous-topos, et les sous-topos sont comptés par des ensembles de flèches au lieu d'ensembles de sommets.

\section{Monoïdes: dénombrements et coins}\label{sec:finite}

Pour un monoïde $M$, le tablier est le topos des $M$-ensembles à droite, $\PSh(\B M)$. La littérature y a déjà appliqué quatre arches et obtenu des réponses complètes: booléen si et seulement si $M$ est un groupe, De Morgan si et seulement si $aM\cap bM\neq\emptyset$ pour tous $a,b$ \cite{hemelaer-rogers}; points essentiels $\leftrightarrow$ idempotents à $\mathcal D$ (ou $\mathcal J$) près \cite{rogers-1905,pirashvili-jncg}; sous-topos $\leftrightarrow$ idéaux bilatères idempotents, pour $M$ fini \cite[Prop.~6.1.2]{pirashvili-jncg}. Ces résultats figurent au journal d'étalonnage (\S\ref{sec:calibration}). Le générateur y ajoute deux choses: les \emph{spécialisations} aux monoïdes classiques, dont nous n'avons pas trouvé le calcul, et un énoncé de transfert.

\subsection{Dénombrer les points essentiels et les sous-topos}

Un idempotent $e$ donne le point essentiel $X\mapsto Xe = \{x : x\cdot e = x\}$, et $e$, $f$ donnent des points isomorphes si et seulement si
\[
e\approx f\quad:\iff\quad\exists\, a,b\in M,\ \ a = fae,\ \ b = ebf,\ \ ba = e,\ \ ab = f,
\]
ce qui, pour des idempotents, est la relation de Green $\mathcal D$ \cite[Prop.~3.1]{costa-steinberg}.

\begin{theorem}\label{th:fibre}
Le nombre de points essentiels, à isomorphisme près, et le nombre de sous-topos des $M$-ensembles à droite sont les suivants:
\begin{center}\small
\begin{tabular}{lll}
\toprule
monoïde $M$ & points essentiels & sous-topos\\
\midrule
$T_n$ (toutes les applications d'un ensemble à $n$ éléments) & $n$ & $n+1$\\
$PT_n$ (applications partielles) & $n+1$ & $n+2$\\
$I_n$ (bijections partielles) & $n+1$ & $n+2$\\
$O_n$ (applications croissantes d'une chaîne à $n$ éléments) & $n$ & $n+1$\\
$M_n(\mathbb F_q)$ pour la multiplication & $n+1$ & $n+2$\\
le monoïde bicyclique $\langle p,q\mid pq=1\rangle$ & $1$ & ---\\
\bottomrule
\end{tabular}
\end{center}
\end{theorem}

\begin{trace}
\item[tablier] $M$-ensembles à droite.
\item[arche] Pont~\ref{br:pts} (points essentiels $\simeq\Kar(\B M)^{\op}$), et pont~\ref{br:sub}.
\item[spécialiser] objets de $\Kar(\B M)$: les idempotents; isomorphisme: $\approx$. Lean: \lean{nonempty_essentialPoint_iso_iff}, \lean{exists_iso_essentialPoint} (\lean{EssentialPoints.lean}).
\item[réécrire] compter les classes d'idempotents pour $\approx$, monoïde par monoïde: ce sont les rangs.
\item[vérifier] Lean, pour les matrices sur un corps quelconque: \lean{karoubiRel_iff_rank_eq} ($e\approx f\iff\operatorname{rk}e=\operatorname{rk}f$) et \lean{essentialPointClassesEquivFin} ($n+1$ classes de points essentiels) (\lean{EssentialPoints.lean}); pour $T_n$, \lean{nonempty_karoubiMk_iso_iff_card_eq} et \lean{karoubiClassesEquivFin} (exactement $n$ classes) (\lean{KaroubiMonoid.lean}); les autres lignes, et tous les dénombrements de sous-topos, par la preuve directe seule.
\item[nouveauté] les correspondances sont publiées \cite{rogers-1905,pirashvili-jncg}, avec l'exemple $T_3$ \cite[Ex.~6.1.5]{pirashvili-jncg}; les dénombrements pour $n$ quelconque et pour ces monoïdes n'ont pas été trouvés.
\end{trace}

\direct{} Avec $ab = a\circ b$: si $e\approx f$, alors $b\circ a = e$ est l'identité sur $\operatorname{im}e$, donc $a$ envoie $\operatorname{im}e$ injectivement dans $\operatorname{im}f$, et symétriquement; les images ont le même cardinal (le rang). Réciproquement, pour une bijection $\theta : \operatorname{im}e\to\operatorname{im}f$ (un isomorphisme d'ordres pour $O_n$, un isomorphisme linéaire pour les matrices), $a = \theta\circ e$ et $b = \theta^{-1}\circ f$ appartiennent au monoïde et vérifient les quatre équations. Il existe des idempotents d'images de cardinal $1,\dots,n$ dans $T_n$, $O_n$, et $0,\dots,n$ dans $PT_n$, $I_n$, $M_n(\mathbb F_q)$. Dans le monoïde bicyclique, les idempotents sont les $q^ip^i$, et $a = q^i$, $b = p^i$ donnent $q^ip^i\approx 1$. Quant aux sous-topos, pour $M$ fini ils correspondent aux idéaux bilatères $I$ tels que $I^2 = I$, auxquels s'ajoute l'idéal vide (le sous-topos dégénéré) \cite[Prop.~6.1.2]{pirashvili-jncg}. Dans chacun de ces monoïdes, les idéaux bilatères non vides sont les idéaux de rang $\{m : \operatorname{rk}m\le k\}$, qui forment une chaîne, et chacun est engendré par un idempotent $e$ de rang $k$, donc idempotent ($I\supseteq I^2\ni e$, d'où $I^2\supseteq MeM = I$). En comptant les rangs possibles et en ajoutant $\emptyset$, on obtient la dernière colonne.
\qed

\subsection{Transfert le long des coins}

\begin{theorem}\label{th:transfer}
Soit $e$ un idempotent d'un monoïde $M$ tel que $1\in MeM$, disons $1 = res$. Alors $aM\cap bM\neq\emptyset$ pour tous $a,b\in M$ si et seulement si $a(eMe)\cap b(eMe)\neq\emptyset$ pour tous $a,b\in eMe$.
\end{theorem}

\begin{trace}
\item[tablier] $\PSh(\B M)\simeq\PSh(\B(eMe))$: avec $a = esf$, $b = fre$, tout idempotent $f$ de $M$ est isomorphe dans $\Kar(\B M)$ à l'idempotent $esfre$ de $eMe$; puis le pont~\ref{br:morita}.
\item[arche] invariance de Morita de la condition d'Ore (\lean{rightOreCondition_iff_of_moritaEquivalentCat}).
\item[vérifier] preuve directe ci-dessous; Lean: \lean{forall_exists_mul_eq_mul_iff_cornerMonoid} (l'énoncé algébrique), et comme corollaire des ponts \lean{karoubiCornerEquiv}, \lean{moritaEquivalentCat_cornerMonoid}, \lean{rightOreCondition_singleObj_iff_cornerMonoid} (\lean{Corners.lean}).
\item[nouveauté] conséquence immédiate de résultats publiés (théorème de Morita pour les monoïdes \cite{knauer,banaschewski,rogers-1905} et \cite[Th.~2.31]{hemelaer-rogers}); non trouvé sous forme d'énoncé.
\end{trace}

\direct{} Si $M$ vérifie la condition et si $a,b\in eMe$, on prend $ac = bd$; alors $x = ecre$, $y = edre\in eMe$ et $ax = acre = bdre = by$ (car $a = ae$, $b = be$). Si $eMe$ vérifie la condition et si $a,b\in M$, on l'applique à $esare$, $esbre$: $esare\,x = esbre\,y$ avec $x,y\in eMe$; on multiplie à gauche par $r$ et on utilise $res = 1$: $a(rex) = b(rey)$.\qed

Dans un monoïde fini, $1\in MeM$ impose $e = 1$; ce résultat concerne donc les monoïdes infinis, par exemple les monoïdes d'endomorphismes d'espaces vectoriels de dimension infinie.

\section{Analyse: faisceaux sur les espaces compacts et treillis de Banach}\label{sec:analysis}

Jusqu'ici, les tabliers étaient des topos de préfaisceaux. En analyse, les tabliers naturels sont des topos de \emph{faisceaux} sur un espace topologique $X$: un faisceau associe un ensemble $F(U)$ à chaque ouvert, avec des restrictions, de sorte que des sections locales compatibles se recollent de façon unique (fonctions continues, solutions d'une équation différentielle, sections d'un fibré). La négation se calcule sur les ouverts: le pseudo-complémentaire d'un ouvert $U$ est l'intérieur de son complémentaire. Deux caractérisations classiques des topos de faisceaux sur un espace jouent ici le rôle d'arches (ce sont les pendants topologiques des ponts~\ref{br:bool} et~\ref{br:dm}; elles ne font pas partie des cinq arches démontrées en Lean et ne sont pas formalisées):
\begin{itemize}
\item $\Sh(X)$ est booléen si et seulement si toute partie ouverte de $X$ est fermée (le pseudo-complémentaire de $U$ est $X\setminus\overline U$, et $U\cup(X\setminus\overline U) = X$ si et seulement si $\overline U\subseteq U$);
\item $\Sh(X)$ est de De Morgan si et seulement si $X$ est \emph{extrêmement discontinu}: l'adhérence de tout ouvert est ouverte \cite{johnstone-demorgan}. Sur les ouverts, $\neg U = X\setminus\overline U$ et $\neg\neg U = \operatorname{int}\overline U$, et $\neg U\cup\neg\neg U = X$ si et seulement si $\overline U\subseteq\operatorname{int}\overline U$, c'est-à-dire si et seulement si $\overline U$ est ouvert.
\end{itemize}
Le second registre est l'analyse fonctionnelle: par la dualité de Gelfand, un espace compact $X$ revient au même que la C$^*$-algèbre commutative $C(X)$, ou que le treillis de Banach $C(X,\mathbb R)$.

\begin{theorem}\label{th:analysis}
Soit $X$ un espace compact (séparé).
\begin{enumerate}
\item $\Sh(X)$ est booléen si et seulement si $X$ est fini, si et seulement si $C(X)$ est de dimension finie.
\item $\Sh(X)$ est de De Morgan si et seulement si le treillis $C(X,\mathbb R)$ est \emph{complet au sens de Dedekind}: toute famille majorée de fonctions continues a une borne supérieure parmi les fonctions continues.
\item Si $X$ est métrisable, $\Sh(X)$ est de De Morgan si et seulement si $X$ est fini. En particulier, les topos des faisceaux sur $[0,1]$, sur le cercle, sur l'ensemble de Cantor ne sont pas de De Morgan, tandis que les topos des faisceaux sur le compactifié de Stone--Čech $\beta\mathbb N$ et sur le spectre de Gelfand de $L^\infty[0,1]$ sont de De Morgan et non booléens.
\end{enumerate}
\end{theorem}

\begin{trace}
\item[tablier] $\Sh(X)$, $X$ compact, avec $X = $ le spectre de Gelfand de $C(X)$.
\item[arche] les deux caractérisations classiques ci-dessus (ouverts fermés; extrêmement discontinu \cite{johnstone-demorgan}).
\item[spécialiser] pour $X$ compact: tout ouvert est fermé $\iff$ $X$ discret $\iff$ $X$ fini; extrêmement discontinu $\iff$ $C(X,\mathbb R)$ complet au sens de Dedekind (\cite[Th.~14]{stone-bounded}, et Nakano \cite{nakano}); un espace compact métrisable extrêmement discontinu est fini.
\item[réécrire] tout énoncer sur $C(X)$.
\item[vérifier] preuve directe ci-dessous; absent de la bibliothèque Lean (Mathlib a \texttt{ExtremallyDisconnected} et \texttt{Stonean}, mais pas $\Sh(X)$ comme topos muni de sa négation).
\item[nouveauté] la forme interne est publiée: Johnstone montre que $\Sh(X)$ est de De Morgan si et seulement si son objet des réels de Dedekind est intérieurement conditionnellement complet \cite{johnstone-demorgan}; la forme externe (2), énoncée sur $C(X,\mathbb R)$, qui combine ce résultat avec ceux de Stone et de Nakano \cite{stone-bounded,nakano}, et le cas métrisable (3) n'ont pas été trouvés.
\end{trace}

\direct{} (1) Si tout ouvert est fermé, tout singleton (fermé, car $X$ est séparé) a un complémentaire ouvert qui est fermé, donc les singletons sont ouverts; $X$ est discret et compact, donc fini. Alors $C(X)\cong\mathbb C^{X}$. (2) C'est un corollaire de deux théorèmes classiques: $\Sh(X)$ est de De Morgan si et seulement si $X$ est extrêmement discontinu (Johnstone \cite{johnstone-demorgan}; le calcul de la négation sur les ouverts est rappelé plus haut), et, pour $X$ compact séparé, $X$ est extrêmement discontinu si et seulement si $C(X,\mathbb R)$ est complet au sens de Dedekind (Stone \cite[Th.~14]{stone-bounded} et Nakano \cite{nakano}). Nous citons le sens \q{extrêmement discontinu implique complet au sens de Dedekind}; l'autre se voit facilement. Pour un ouvert $U$, tout majorant continu des fonctions $0\le g\le 1$ à support dans $U$ est $\ge 1$ sur $\overline U$ (Urysohn), et, pour $x\notin\overline U$, l'un de ces majorants s'annule en $x$; une borne supérieure dans $C(X,\mathbb R)$ est donc nécessairement la fonction indicatrice de $\overline U$, qui n'est continue que si $\overline U$ est ouvert. (3) Dans un espace métrisable extrêmement discontinu, toute suite convergente de points distincts $x_n\to x$ donne des ouverts disjoints $U$, $V$ autour de $\{x_{2n}\}$ et de $\{x_{2n+1}\}$ dont les adhérences contiennent toutes deux $x$; or, dans un tel espace, deux ouverts disjoints ont des adhérences disjointes. Un espace compact métrisable de ce type n'a donc aucune suite convergente non triviale, et il est fini. L'espace $\beta\mathbb N$ et le spectre de $L^\infty[0,1]$ sont extrêmement discontinus (le second parce que $L^\infty[0,1]$ est complet au sens de Dedekind), et infinis.\qed

\inwords{} La logique intuitionniste des \q{fonctions continues sur $X$} vérifie $\neg p\vee\neg\neg p$ exactement quand les bornes supérieures de fonctions continues restent continues. Sur $[0,1]$, c'est faux: les fonctions $\min(1, n\cdot\max(0, x-\tfrac12))$ sont majorées par $1$, et leur borne supérieure, l'indicatrice de $(\tfrac12,1]$, n'est pas continue; sur le spectre de $L^\infty$, où toute famille bornée a une borne supérieure essentielle, c'est vrai.

\subsection{Un espace, trois théorèmes classiques}

Il s'agit d'une équivalence classique du type discuté au \S\ref{sec:equivalences}. Un espace sobre est déterminé par son topos des faisceaux: $\Sh(X)\simeq\Sh(Y)$ si et seulement si $X\cong Y$, les points de $\Sh(X)$ étant les points de $X$. Pour les espaces compacts (qui sont sobres), on lit le même invariant, l'espace des points, sur trois présentations:
\begin{itemize}
\item sur $\Sh(X)$: deux espaces compacts dont les topos des faisceaux sont équivalents sont homéomorphes;
\item sur l'anneau $C(X)$: le \emph{théorème de Gelfand--Kolmogorov}, $C(X)\cong C(Y)$ comme anneaux si et seulement si $X\cong Y$ (les points sont les idéaux maximaux);
\item sur l'espace de Banach $C(X)$: le \emph{théorème de Banach--Stone}, $C(X)$ et $C(Y)$ sont linéairement isométriques si et seulement si $X\cong Y$ (les points sont les points extrémaux de la boule unité du dual, au signe près).
\end{itemize}
Les trois énoncés sont classiques et démontrés dans trois cours différents; le pont dit qu'ils sont un seul énoncé, \q{l'espace des points est un invariant}, lu sur trois présentations du même objet.

\section{Algèbre linéaire et arithmétique}\label{sec:linear}

Les topos sont faits d'ensembles; l'algèbre linéaire y entre donc par des structures \emph{ensemblistes} attachées aux espaces vectoriels: actions de monoïdes de matrices, foncteurs sur des treillis de sous-espaces. Les réponses diffèrent des réponses linéaires.

\subsection{Monoïdes de matrices: l'équivalence de Morita garde la mémoire de la taille}

Pour les anneaux, la théorie de Morita dit que les $M_n(k)$-modules et les $k$-modules forment des catégories équivalentes, pour tout $n$. Pour les ensembles, c'est le contraire.

\begin{theorem}\label{th:matrixmorita}
Soit $k$ un corps et $M_n(k)$ le monoïde des matrices $n\times n$ pour la multiplication.
\begin{enumerate}
\item Les $M_n(k)$-ensembles à droite sont équivalents aux préfaisceaux sur la catégorie $\mathsf{Vect}_{\le n}$ des $k$-espaces vectoriels de dimension au plus $n$.
\item Les $M_n(k)$-ensembles à droite et les $M_m(k)$-ensembles à droite forment des catégories équivalentes si et seulement si $n = m$.
\item Les $\mathrm{GL}_n(k)$-ensembles à droite forment un topos booléen dont le seul point essentiel (à isomorphisme près) est le foncteur d'oubli, de groupe d'automorphismes $\mathrm{GL}_n(k)$.
\end{enumerate}
\end{theorem}

\begin{trace}
\item[tablier] $\PSh(\B M_n(k))$.
\item[arche] Pont~\ref{br:morita} (enveloppe de Karoubi), pont~\ref{br:pts}, pont~\ref{br:bool}.
\item[spécialiser] $\Kar(\B M_n(k))$: les objets sont les projecteurs $e$, les flèches $e\to f$ sont les matrices $a$ telles que $a = fae$, c'est-à-dire les applications linéaires $\operatorname{im}e\to\operatorname{im}f$; donc $e\mapsto\operatorname{im}e$ est une équivalence $\Kar(\B M_n(k))\simeq\mathsf{Vect}_{\le n}$. Classes d'isomorphisme: les dimensions $0,\dots,n$.
\item[réécrire] (2): des enveloppes de Karoubi équivalentes auraient le même nombre, $n+1$ et $m+1$, de classes d'isomorphisme d'objets.
\item[vérifier] Lean: \lean{karoubiRel_iff_rank_eq} (deux projecteurs sont isomorphes dans $\Kar$ si et seulement s'ils ont le même rang), \lean{essentialPointEndMulEquivOfGroup}, \lean{isBooleanTopos_singleObj_of_group}; l'équivalence avec $\mathsf{Vect}_{\le n}$ par l'argument direct.
\item[nouveauté] non trouvé sous cette forme; (2) découle du théorème de Morita pour les monoïdes, puisque dans $M_n(k)$ un inverse d'un côté est un inverse, de sorte que l'équivalence de Morita est l'isomorphisme \cite[Cor.~1.6.2]{rogers-thesis}. L'analogue linéaire (l'enveloppe de Karoubi des $k[M_n(\mathbb F_q)]$-modules) est étudié dans \cite{kuhn}; la description ensembliste (1) n'a pas été trouvée.
\end{trace}

\direct{} Une flèche $e\to f$ dans $\Kar(\B M_n(k))$ est une matrice $a$ telle que $fae = a$; sa restriction à $\operatorname{im}e$ est une application linéaire $\operatorname{im}e\to\operatorname{im}f$, et toute application linéaire provient d'un unique tel $a$ (on prolonge par $0$ sur $\ker e$). Tout espace de dimension $\le n$ est l'image d'un projecteur. Donc $\Kar(\B M_n(k))\simeq\mathsf{Vect}_{\le n}$, et le pont~\ref{br:morita} donne (1). Pour (2), une équivalence de catégories préserve le nombre de classes d'isomorphisme d'objets, $n+1$ d'un côté et $m+1$ de l'autre. (3) est le pont~\ref{br:bool} et le pont~\ref{br:pts} pour un groupe.\qed

\inwords{} La théorie des représentations ensemblistes du monoïde de matrices voit la taille des matrices; la théorie linéaire ne la voit pas. La différence vient de ce que les idempotents se scindent dans $\mathsf{Vect}$, mais que la catégorie des ensembles n'\q{additionne} pas les objets: $k^n\oplus k^n\cong k^{2n}$ est invisible pour les $M_n(k)$-ensembles.

\subsection{Treillis de sous-espaces}

Soit $\mathrm{Sub}(V)$ le treillis des sous-espaces d'un espace vectoriel $V$ de dimension finie, et considérons le topos $[\mathrm{Sub}(V),\Set]$ des foncteurs sur ce treillis, c'est-à-dire des faisceaux sur son espace d'Alexandrov.

\begin{theorem}\label{th:subspaces}
\begin{enumerate}
\item $[\mathrm{Sub}(V),\Set]$ est de De Morgan et, si $V\neq 0$, non booléen.
\item $[\mathrm{Sub}(V),\Set]\simeq[\mathrm{Sub}(W),\Set]$ si et seulement si $\mathrm{Sub}(V)\cong\mathrm{Sub}(W)$; le topos détermine $\dim V$ (la hauteur du treillis) et, si $\dim V\ge 3$, le corps à isomorphisme près; en dimension $2$ il ne détermine que $|k|$, donc $k$ lorsque $k$ est fini.
\item Ses points essentiels sont les évaluations aux sous-espaces de $V$; sur $\mathbb F_q$, il y en a $\sum_{k=0}^{n}\binom nk_q$ (coefficients binomiaux de Gauss).
\end{enumerate}
\end{theorem}

\begin{trace}
\item[tablier] $[\mathrm{Sub}(V),\Set] = \PSh(\mathrm{Sub}(V)^{\op})$.
\item[arche] Pont~\ref{br:dm} (un treillis est confluent), pont~\ref{br:bool}, pont~\ref{br:morita} (un ensemble ordonné est sa propre enveloppe de Karoubi), pont~\ref{br:pts}.
\item[réécrire] (2) par le théorème fondamental de la géométrie projective: pour $\dim V\ge 3$, un isomorphisme de treillis $\mathrm{Sub}(V)\cong\mathrm{Sub}(W)$ est induit par une bijection semi-linéaire \cite{artin-geometric}; sur $\mathbb F_q$ avec $n = 2$, le nombre $q+1$ de droites détermine $q$.
\item[vérifier] Lean: \lean{isDeMorganTopos_functor_preorder_iff}, \lean{moritaEquivalentCat_preorder_iff}; le reste par l'argument direct.
\item[nouveauté] non trouvé; les ingrédients sont classiques (un ensemble ordonné est sa propre enveloppe de Karoubi; théorème fondamental de la géométrie projective \cite{artin-geometric}).
\end{trace}

\subsection{Le monoïde multiplicatif de $\mathbb Z/n$}

\begin{theorem}\label{th:zn}
Soit $n\ge 2$ et soit $\omega(n)$ le nombre de facteurs premiers distincts de $n$. Le topos des ensembles munis d'une action du monoïde multiplicatif $(\mathbb Z/n,\cdot)$
\begin{enumerate}
\item est de De Morgan et non booléen;
\item a exactement $2^{\omega(n)}$ points essentiels à isomorphisme près;
\item a exactement $D(\omega(n))$ sous-topos, où $D(k)$ est le $k$-ième nombre de Dedekind, le nombre de fonctions booléennes monotones de $k$ variables: $D(1) = 3$, $D(2) = 6$, $D(3) = 20$, $D(4) = 168$, $D(5) = 7581$.
\end{enumerate}
\end{theorem}

\begin{trace}
\item[tablier] $\PSh(\B(\mathbb Z/n,\cdot))$.
\item[arche] Ponts~\ref{br:dm}, \ref{br:bool}, \ref{br:pts} et, pour les monoïdes finis, la classification des sous-topos par les idéaux bilatères idempotents \cite[Prop.~6.1.2]{pirashvili-jncg}.
\item[spécialiser] $0$ est un zéro, donc Ore est vérifiée; la classe d'un facteur premier de $n$ n'est pas inversible. Le monoïde est commutatif, donc les points essentiels sont les idempotents; par le théorème chinois, les idempotents de $\mathbb Z/n$ correspondent aux parties de l'ensemble des facteurs premiers de $n$. Idéaux idempotents: voir ci-dessous.
\item[vérifier] preuve directe ci-dessous; Lean: non formalisé (la classification des sous-topos des monoïdes finis par les idéaux idempotents n'est pas dans la bibliothèque); les dénombrements sont vérifiés par ordinateur.
\item[nouveauté] la théorie générale est publiée (\cite[Prop.~6.1.2, Prop.~6.3.1, Th.~5.6.1]{pirashvili-jncg}: sous-topos $\leftrightarrow$ parties commençantes de l'ensemble ordonné des $\mathcal J$-classes d'idempotents); le monoïde $(\mathbb Z/n,\cdot)$, les décomptes et les nombres de Dedekind n'ont pas été trouvés. Vérifié par ordinateur pour $n\le 60$ et $n = 210, 2310$, et, directement à partir des axiomes d'une topologie de Grothendieck, pour $n\le 30$.
\end{trace}

\direct{} (1) et (2): comme dans la trace. Pour (3), les sous-topos correspondent aux idéaux bilatères $I$ du monoïde fini $M = (\mathbb Z/n,\cdot)$ tels que $I^2 = I$, l'idéal vide compris \cite[Prop.~6.1.2]{pirashvili-jncg}. Dans un monoïde fini, un idéal idempotent est engendré par ses idempotents: si $I = I^2$, tout $x\in I$ est un produit d'éléments de $I$ en nombre arbitrairement grand, et un produit assez long dans un demi-groupe fini se factorise en $uev$ avec $e$ idempotent, $e$ produit d'éléments de $I$. Les idéaux idempotents sont donc les réunions d'idéaux principaux $eM$, $e$ idempotent, et $eM\subseteq fM$ si et seulement si $e = ef$. En écrivant $e\leftrightarrow S_e\subseteq\{\text{facteurs premiers de }n\}$ (les premiers $p$ tels que $e\equiv 1\bmod p^{v_p(n)}$), on a $e = ef$ si et seulement si $S_e\subseteq S_f$; les idéaux idempotents correspondent donc aux familles, stables vers le bas, de parties d'un ensemble à $\omega(n)$ éléments, dont le nombre est le nombre de Dedekind $D(\omega(n))$.\qed

\begin{example}
Pour $n = p$ premier: trois sous-topos (tous les $(\mathbb Z/p,\cdot)$-ensembles; l'idéal $\{0\}$; le sous-topos dégénéré). Pour $n = p^2$, trois aussi: l'idéal $p\,\mathbb Z/p^2$ n'est pas idempotent, puisque son carré est $\{0\}$. Pour $n = 30$: $2^3 = 8$ points essentiels et $D(3) = 20$ sous-topos.
\end{example}

\subsection{Autres domaines du programme de troisième année}

La même méthode s'applique partout où une structure familière est un topos ayant une seconde présentation. Une carte sommaire du programme de troisième année de licence:

\begin{center}\small
\begin{tabularx}{\linewidth}{>{\raggedright\arraybackslash}p{3.6cm}>{\raggedright\arraybackslash}X}
\toprule
domaine & tabliers\\
\midrule
théorie des groupes & $G$-ensembles, Cayley, décomposition en orbites (\S\ref{sec:equivalences})\\
anneaux, corps, Galois & monoïdes multiplicatifs, catégories galoisiennes (\S\ref{sec:equivalences})\\
algèbre linéaire & monoïdes de matrices, treillis de sous-espaces (cette section)\\
arithmétique & $(\mathbb Z/n,\cdot)$, ordres de divisibilité (cette section, \S\ref{sec:calibration})\\
topologie générale & faisceaux sur des espaces, espaces d'Alexandrov (\S\ref{sec:analysis}, \S\ref{sec:calibration})\\
analyse fonctionnelle & $C(X)$ et dualité de Gelfand (\S\ref{sec:analysis})\\
mesure et probabilités & algèbres de mesure, qui sont des algèbres de Boole complètes, et leurs topos booléens de faisceaux\\
équations différentielles & faisceaux de solutions; systèmes dynamiques discrets (\S\ref{sec:machines})\\
combinatoire, graphes & carquois, graphes, ensembles simpliciaux (\S\ref{sec:machines}, \S\ref{sec:shapes})\\
logique & topos classifiants, quotients et topologies (\S\ref{sec:calibration})\\
\bottomrule
\end{tabularx}
\end{center}

\section{Une douzaine de topos combinatoires}\label{sec:shapes}

Le dernier lot montre le générateur sous son jour le plus mécanique : une forme venue de la combinatoire, la condition d'Ore vérifiée sur elle, le tiers exclu faible lu directement. Nous n'avons trouvé aucun de ces exemples dans la littérature, sauf les graphes réflexifs \cite[Ex.~2.52]{hemelaer-rogers}.

\begin{theorem}\label{th:shapes}
\begin{center}\small
\begin{tabular}{llcc}
\toprule
préfaisceaux & forme & booléen & De Morgan\\
\midrule
graphes orientés & $V\rightrightarrows E$ & non & non\\
graphes réflexifs & $V\rightrightarrows E$ avec boucles & non & non\\
ensembles simpliciaux & $\Delta$ & non & non\\
ensembles semi-simpliciaux & $\Delta_{\mathrm{inj}}$ & non & non\\
ensembles simpliciaux symétriques & ensembles finis non vides & non & non\\
ensembles cubiques & catégorie des cubes & non & non\\
ensembles bisimpliciaux & $\Delta\times\Delta$ & non & non\\
ensembles simpliciaux augmentés & $\Delta_+$ & non & oui\\
$\mathsf{FI}$-ensembles & ensembles finis et injections & non & oui\\
préfaisceaux sur les ensembles finis & $\mathsf{Fin}$ & non & oui\\
le topos des arbres & $\omega$ & non & oui\\
ensembles munis d'une involution & $\B(\mathbb Z/2)$ & oui & oui\\
\bottomrule
\end{tabular}
\end{center}
\end{theorem}

\begin{trace}
\item[tablier] chaque forme.
\item[arche] Ponts~\ref{br:bool} et~\ref{br:dm}.
\item[spécialiser] \q{non} : deux flèches distinctes issues d'un sommet vers un même objet (source et but d'une arête, les extrémités $[0]\rightrightarrows[1]$) qu'aucune flèche n'égalise, toute flèche vers un sommet étant constante. \q{oui} : un objet initial (on complète toute paire de flèches de même but à partir de lui), ou une chaîne (on complète par le minimum).
\item[vérifier] Lean : \lean{not_isDeMorganTopos_graphs} avec le contre-exemple explicite \lean{graphs_pseudocompl_counterexample}, \lean{not_isDeMorganTopos_simplicialSets}, \lean{isDeMorganTopos_augmentedSimplicialSets}, \lean{isDeMorganTopos_presheaf_fintypeCatSkeleton}, \lean{isDeMorganTopos_treeTopos}, \lean{not_isBooleanTopos_treeTopos}, tous par le résultat général \lean{isDeMorganTopos_presheaf_of_hasInitial} ; les autres lignes par la preuve directe.
\item[nouveauté] le critère est classique ; les exemples, sauf les graphes réflexifs, n'ont pas été trouvés.
\end{trace}

\direct{} Pour un \q{non}, dans le graphe à une arête (ou dans $\Delta[1]$), on prend pour $A$ une extrémité $u$ : alors $\neg A$ est l'autre extrémité $v$, $\neg\neg A = \{u\}$, et l'arête n'est dans aucun des deux. Pour un \q{oui} avec un objet initial $\emptyset$, tout élément $x$ a une restriction $\res x\emptyset$ qui est aussi une restriction de toute restriction de $x$ ; si une restriction de $x$ est dans $A$, alors $\res x\emptyset$ l'est aussi, et toute restriction de $x$ admet alors une restriction ultérieure dans $A$. Pour $\omega$, deux restrictions de $x$ ont une restriction ultérieure commune au niveau inférieur.\qed

\begin{example}
Dans les ensembles simpliciaux augmentés, un simplexe appartient à $\neg\neg A$ si et seulement si son étiquette d'augmentation appartient à $A_{-1}$, et à $\neg A$ sinon : le tiers exclu faible est décidé par l'augmentation. Dans le topos des arbres, la logique interne démontre $\neg\varphi\vee\neg\neg\varphi$ pour tout $\varphi$, mais pas $\varphi\vee\neg\varphi$.
\end{example}

\section{Des théorèmes classiques comme deux lectures d'un même pont}\label{sec:equivalences}

La même machinerie explique pourquoi certaines paires de théorèmes de troisième année de licence, enseignés dans des cours différents, sont \q{le même théorème}: ce sont deux lectures d'un même invariant sur un même topos. Rien dans cette section n'est nouveau; son but est pédagogique, et elle montre ce que fait un pont quand ses deux extrémités sont déjà connues.

\subsection{Relations d'équivalence et quotients}

Soit $E$ une relation d'équivalence sur un ensemble $S$, vue comme un groupoïde (une flèche $s\to t$ quand $s\mathrel E t$). Un préfaisceau sur $E$ est une famille d'ensembles $(F_s)_{s\in S}$ munie de bijections compatibles $F_s\cong F_t$ pour $s\mathrel E t$.

\begin{itemize}
\item \emph{Tablier:} le groupoïde $E$ est équivalent à la catégorie discrète $S/E$ (chaque composante connexe est contractile), donc, d'après le pont~\ref{br:morita}, $\PSh(E)\simeq\Set^{S/E}$.
\item \emph{Lecture 1 (algèbre de L1--L2):} les familles sur $S$ compatibles avec $E$ sont la même chose que les familles sur $S/E$; pour $F_s = \{\ast\}$ ou $\emptyset$, c'est la propriété universelle du quotient: \emph{une application $S\to\{0,1\}$ constante sur les classes se factorise par $S/E$}.
\item \emph{Lecture 2 (pont~\ref{br:bool}):} $E$ est un groupoïde, donc tout sous-préfaisceau a un complémentaire; pour les parties, \emph{une réunion de classes a pour complémentaire une réunion de classes}: les classes forment une partition de $S$.
\end{itemize}

\subsection{Orbites et complémentaires}

Pour un groupe $G$, le pont~\ref{br:bool} dit que le complémentaire d'un sous-$G$-ensemble est un sous-$G$-ensemble. Avec la décomposition de tout préfaisceau en composantes connexes (vraie dans tout topos de préfaisceaux), c'est la \emph{décomposition en orbites}: tout $G$-ensemble est la réunion disjointe de ses orbites, et tout sous-$G$-ensemble est une réunion d'orbites. Pour un monoïde qui n'est pas un groupe, la décomposition en composantes connexes reste vraie, mais une composante n'est pas forcément \q{indécomposable} au sens booléen: le pont montre que le théorème de décomposition en orbites est un énoncé sur l'inversibilité.

\subsection{Le théorème de Birkhoff et les espaces $T_0$ finis}

Soit $P$ un ensemble ordonné fini et $X_P$ son espace d'Alexandrov (ouverts $=$ parties stables vers le haut, ou sections finissantes). Le tablier est $\Sh(X_P)\simeq[P,\Set]$: un faisceau sur $X_P$ est déterminé par ses valeurs sur les plus petits voisinages ouverts $\uparrow p$, et $\uparrow q\subseteq\uparrow p$ quand $p\le q$. Deux invariants de ce topos:
\begin{itemize}
\item ses \emph{objets sous-terminaux} (sous-préfaisceaux du préfaisceau à un point): lus sur $X_P$, ce sont les ouverts; lus sur $[P,\Set]$, ce sont les parties de $P$ stables vers le haut;
\item ses \emph{points}: lus sur $X_P$ (un espace sobre), ce sont les points de $X_P$; lus sur $[P,\Set]$, ce sont les évaluations aux éléments de $P$, et les parties principales $\uparrow p$ sont exactement les éléments sup-irréductibles de $\mathrm{Up}(P)$, avec $\uparrow p\subseteq\uparrow q$ si et seulement si $q\le p$.
\end{itemize}
Lire ces invariants à rebours donne d'un coup deux théorèmes classiques: le \emph{théorème de représentation de Birkhoff} (un treillis distributif fini $L$ est isomorphe au treillis des parties stables vers le haut de l'opposé de l'ensemble ordonné de ses éléments sup-irréductibles) et la \emph{correspondance entre espaces $T_0$ finis et ensembles ordonnés finis} (un espace $T_0$ fini est l'espace d'Alexandrov de son ordre de spécialisation). En appliquant de plus le pont~\ref{br:bool}: le topos est booléen si et seulement si $P$ est discret, si et seulement si $\mathrm{Up}(P)$ est l'ensemble des parties de $P$, ce qui est le cas fini du \emph{théorème de représentation de Stone}: une algèbre de Boole finie est l'ensemble des parties de ses atomes.

\paragraph{Un précédent dans les notes de Grothendieck.} Dans une note manuscrite non datée sur les topos \cite{grothendieck-161-3}, p.~26, Grothendieck énumère des conditions équivalentes sur un topos $E$: être équivalent à la fois au topos d'un espace et à un topos de préfaisceaux; être le topos d'un espace \q{essentiel} (un espace dont tous les points sont essentiels); être équivalent à $\widehat I$ pour un ensemble ordonné $I$; être engendré par les sous-objets de l'objet final et avoir suffisamment de points essentiels. Il demande de prouver que $\widehat C$ satisfait ces conditions si et seulement si $C$ est préordonnée, et que le topos d'un espace $X$ les satisfait si et seulement si les points essentiels de $X$ sont très denses, et il note des 2-équivalences entre topos spatiaux essentiels, préordres et espaces essentiels. Ce sont des ponts au sens de cet article, lus sur le tablier des préordres avec les arches des sous-objets de l'objet final et des points essentiels. Le générateur retrouve le premier exercice par l'arche de Morita: si $\widehat C\simeq\widehat P$ pour un préordre $P$, alors $\Kar(C)\simeq\Kar(P)\simeq P$, et $C$, sous-catégorie pleine de son enveloppe de Karoubi, est mince; réciproquement une catégorie mince est équivalente à son squelette mince, qui est un ordre (Lean: \lean{exists_preorder_moritaEquivalentCat_iff_isThin}). Le second exercice, sur les espaces, n'est pas formalisé ici.

\subsection{Le théorème de Cayley et la reconstruction d'un groupe}

Pour un groupe $G$, le pont~\ref{br:pts} dit deux choses: le foncteur d'oubli $U$ des $G$-ensembles vers les ensembles est, à isomorphisme près, l'unique point essentiel; son groupe d'automorphismes est $G$ (Lean: \lean{essentialPointEndMulEquivOfGroup}). En évaluant les automorphismes de $U$ sur le $G$-ensemble régulier $G$, on obtient un morphisme injectif $G\to\mathrm{Sym}(G)$: c'est le \emph{théorème de Cayley}. Réciproquement, le théorème de Cayley et le lemme de Yoneda (une transformation naturelle de $U = \Hom(G,-)$ est déterminée par sa valeur sur l'identité) donnent la reconstruction. Les deux sont un seul énoncé sur le topos $\PSh(\B G)$, lu sur ses points et sur son objet représentable.

\subsection{Correspondance de Galois et revêtements}

Deux théorèmes de troisième année de licence, le \emph{théorème fondamental de la théorie de Galois} (corps intermédiaires d'une extension galoisienne finie $L/K$ $\leftrightarrow$ sous-groupes de $\mathrm{Gal}(L/K)$) et la \emph{classification des revêtements} (revêtements connexes d'un espace connexe raisonnable $X$ $\leftrightarrow$ classes de conjugaison de sous-groupes de $\pi_1(X)$), sont deux lectures d'un même pont, explicité par Grothendieck \cite{sga1}: la catégorie des $K$-algèbres étales finies déployées par $L$ et la catégorie des revêtements de $X$ sont toutes deux équivalentes à une catégorie de $G$-ensembles ($G = \mathrm{Gal}(L/K)$, resp.\ $\pi_1(X)$); l'invariant est l'ensemble ordonné des \emph{objets connexes pointés}, qui, du côté des $G$-ensembles, est l'ensemble ordonné des sous-groupes. (Ces tabliers ne sont pas des topos de préfaisceaux sur une forme finie, et ce pont n'est pas dans la bibliothèque Lean; voir \cite{szamuely} pour un exposé de niveau licence.)

\subsection{Confluence et espaces extrêmement discontinus}

La dernière paire figure au journal d'étalonnage: un préordre est confluent si et seulement si son espace d'Alexandrov est extrêmement discontinu \cite[Prop.~7.10]{caramello-stone} (Lean: \lean{extremallyDisconnected_iff_confluent}). Lu sur un système de réécriture, le théorème de Church--Rosser pour la $\beta$-réduction devient le fait qu'un espace de $\lambda$-termes est extrêmement discontinu. Ici, le pont fait plus que redémontrer deux théorèmes connus: une propriété issue de la théorie des langages de programmation et une propriété de topologie générale se révèlent être la même propriété.

\begin{remark}
Dans chaque paire, une lecture est en général démontrée dans un cours et l'autre dans un autre, par des arguments sans rapport. Le pont ne donne pas de preuves plus courtes; il explique \emph{pourquoi} les deux énoncés ont la même forme, et il donne une recette pour trouver d'autres paires de ce genre: prendre un topos ayant deux présentations familières, et lire un invariant sur les deux.
\end{remark}

\section{Une sixième arche, tirée des notes de Grothendieck : recoller deux sous-topos}\label{sec:gluing}

Les cinq arches du \S\ref{sec:toolbox} viennent de la littérature. Une sixième vient des notes de Grothendieck sur l'algèbre universelle (dossier 161-2, p.~48), où il se demande quand deux sous-topos $V$, $W$ d'un topos $B$ peuvent être \q{recollés} : quand les morphismes qui sont \emph{localement} dans $V$ ou dans $W$ forment-ils un sous-topos ? La note associée \cite{bridges-notes} (section \q{Anneaux de caractéristique définie}) rend son critère précis et le démontre ; dans le cas où il l'appliquait, les anneaux de caractéristique définie, il est en défaut, ce qui est un cas particulier du fait que la caractéristique~$0$ n'est pas géométriquement définissable \cite[\S8.2]{caramello-flat}. Nous faisons ici tourner ce critère comme une arche.

\begin{bridge}[Recollement]\label{br:glue}
Pour des sous-topos $V$, $W$ d'un topos $B$, les morphismes géométriques $f\colon\mathcal E\to B$ tels que $\mathcal E$ soit recouvert par deux ouverts sur lesquels $f$ se factorise respectivement par $V$ et par $W$ sont les morphismes qui se factorisent par un sous-topos si et seulement s'il existe des sous-objets $U'$, $U''$ de l'objet final de $B$ tels que
\[
W\wedge o(U')\le V,\qquad V\wedge o(U'')\le W,\qquad V\vee W\le o(U'\vee U'').
\]
Ce sous-topos est alors $V\vee W$. Lean : le cœur ordinal, dans un treillis distributif quelconque, est \lean{Grothendieck.Folder161_2.recollement_iff} (dépôt GrothendieckArchivesProofs) ; le passage aux sous-topos utilise la distributivité du treillis des sous-topos et la localité de la factorisation par un sous-topos \cite[\S A4.5]{elephant}, qui ne sont pas formalisées.
\end{bridge}

\subsection{Le tablier : sous-topos donnés par des sous-catégories pleines}

Soient $C$ une petite catégorie où les idempotents se scindent, et $D\subseteq C$ une sous-catégorie pleine stable par isomorphismes et par rétracts. La restriction le long de $D\subseteq C$, avec l'extension de Kan à droite comme adjoint à droite, fait de $\PSh(D)$ un sous-topos de $\PSh(C)$ ; ce sont les sous-topos \emph{essentiels}, et $\PSh(E)\le\PSh(D)$ si et seulement si $E\subseteq D$ \cite{kelly-lawvere}. Un sous-objet de l'objet final de $\PSh(C)$ est un ensemble $S$ d'objets stable par les flèches entrantes, $o(S)\simeq\PSh(S)$, et le sous-topos ouvert de $\PSh(D)$ découpé par $S$ est $\PSh(D\cap S)$.

\begin{theorem}\label{th:cospan}
Soient $D_1$, $D_2$ des sous-catégories pleines de $C$ stables par isomorphismes et par rétracts. Les morphismes qui sont localement dans $\PSh(D_1)$ ou dans $\PSh(D_2)$ forment un sous-topos si et seulement s'il n'existe pas de \emph{cospan mixte}
\[
a\longrightarrow x\longleftarrow b,\qquad a\in D_1\smallsetminus D_2,\quad b\in D_2\smallsetminus D_1,\quad x\in D_1\cup D_2.
\]
\end{theorem}

\begin{trace}
\item[tablier] $\PSh(C)$ avec les sous-topos essentiels $\PSh(D_1)$, $\PSh(D_2)$.
\item[arche] Pont~\ref{br:glue}.
\item[spécialiser] $W\wedge o(S')\le V$ s'écrit $D_2\cap S'\subseteq D_1$, c'est-à-dire que $S'$ évite $D_2\smallsetminus D_1$ ; le plus grand tel $S'$ est l'ensemble des objets qui ne reçoivent aucune flèche de $D_2\smallsetminus D_1$. De même pour $S''$.
\item[réécrire] $V\vee W\le o(S'\cup S'')$ s'écrit $D_1\cup D_2\subseteq S'\cup S''$ : aucun objet de $D_1\cup D_2$ ne reçoit à la fois une flèche de $D_1\smallsetminus D_2$ et une flèche de $D_2\smallsetminus D_1$.
\item[vérifier] argument direct ci-dessous ; Lean : le cœur de l'arche seulement.
\item[nouveauté] l'arche elle-même n'a pas été trouvée imprimée (note associée) ; les corollaires sont évalués un par un.
\end{trace}

\direct{} D'après le pont~\ref{br:glue}, la condition est l'existence de $S'$, $S''$ comme dans la trace, et les trois conditions sont monotones en $S'$, $S''$, de sorte qu'on peut prendre les plus grands choix. Un cospan mixte $a\to x\leftarrow b$ avec $x\in D_1\cup D_2$ place $x$ hors des deux plus grands choix ; réciproquement, un objet de $D_1\cup D_2$ hors des deux reçoit une flèche de chacune des deux différences.\qed

\subsection{Trois lectures}

\begin{corollary}[Préordres et espaces d'Alexandrov]\label{co:alex}
Soient $P$ un préordre, $X_P$ son espace d'Alexandrov (les ouverts sont les sections commençantes), et $A$, $B\subseteq P$. Les faisceaux sur $X_P$ qui sont localement des faisceaux sur le sous-espace $A$ ou sur le sous-espace $B$ forment un sous-topos si et seulement si $A\smallsetminus B$ et $B\smallsetminus A$ sont \emph{séparés dans $A\cup B$} : aucun ne rencontre l'adhérence de l'autre dans $A\cup B$.
\end{corollary}

\noindent Un cospan mixte $a\le x\ge b$ dit que $x$ est dans les adhérences (les sections finissantes) de $a$ et de $b$. \status{} Non : c'est le fait classique que $A\cup B$ est réunion de deux ouverts relatifs contenus dans $A$ et dans $B$ si et seulement si $A\smallsetminus B$ et $B\smallsetminus A$ sont séparés dans $A\cup B$. Le pont lit une condition de recollement de sous-topos comme un axiome de séparation.

\begin{corollary}[Carquois]\label{co:quivgl}
Soient $Q$ un carquois, $C$ la catégorie libre sur $Q$ (les préfaisceaux sur $C$ sont les représentations de $Q^{\op}$), et $D_1$, $D_2$ des ensembles de sommets. La disjonction locale de $\PSh(D_1)$ et $\PSh(D_2)$ est un sous-topos si et seulement si aucun sommet de $D_1\cup D_2$ n'est l'extrémité à la fois d'un chemin partant de $D_1\smallsetminus D_2$ et d'un chemin partant de $D_2\smallsetminus D_1$.
\end{corollary}

\noindent Une catégorie libre n'a pas d'idempotent non trivial, donc toute sous-catégorie pleine est stable par rétracts, et les flèches sont des chemins. La condition se vérifie en temps linéaire par deux parcours du graphe. \status{} Non trouvé, mais sa preuve tient en une ligne une fois le pont~\ref{br:glue} admis.

\begin{corollary}[Topos à deux valeurs et monoïdes]\label{co:twovalued}
Dans un topos dont les seuls sous-objets de l'objet final sont $0$ et $1$, la disjonction locale de deux sous-topos $V$, $W$ est un sous-topos si et seulement si $V\le W$ ou $W\le V$. En particulier, pour un monoïde $M$, la disjonction locale de deux sous-topos essentiels du topos des $M$-ensembles (donnés par des ensembles d'idempotents stables par rétracts dans $\Kar(\B M)$) n'est un sous-topos que si l'un contient l'autre.
\end{corollary}

\noindent Avec seulement deux ouverts, la troisième condition du pont~\ref{br:glue} force $U'$ ou $U''$ à valoir $1$, d'où $W\le V$ ou $V\le W$. Dans $\Kar(\B M)$ il y a une flèche $fe\colon e\to f$ entre deux idempotents quelconques, de sorte que le théorème~\ref{th:cospan} donne la même réponse. \status{} Non ; c'est presque une tautologie.

\begin{remark}
Cette arche a été ajoutée après le reste de l'article et c'est la plus faible : de ses trois sorties, l'une est de la topologie classique, une autre presque tautologique, et la dernière, le corollaire~\ref{co:quivgl}, est un court critère combinatoire dont la preuve est l'arche elle-même. Elle confirme une limite de la méthode discutée au \S\ref{sec:limits} : une arche dont la condition est déjà élémentaire produit des théorèmes élémentaires.
\end{remark}

\section{Une septième arche, tirée de la théorie de l'homotopie : les catégories test}\label{sec:homotopical}

Les \emph{catégories test} de Grothendieck sont les petites catégories $A$ dont les préfaisceaux modélisent les types d'homotopie ; $\Delta$ est le premier exemple. Leur caractérisation lit un objet de $\PSh(A)$, le classifiant des sous-objets, de façon homotopique ; on peut donc la faire tourner comme une arche dont l'autre côté est la lecture logique du même objet que donne le pont~\ref{br:dm}. La note associée \cite[section \q{Catégories test et arche De Morgan}]{bridges-notes} le fait sur les questions des dossiers 108, 118 et 134-2 de Grothendieck ; nous en donnons ici la sortie.

Fixons un localisateur fondamental $W$ sur les petites catégories, autre que les deux localisateurs grossiers (par exemple les foncteurs dont le nerf est une équivalence faible). Pour un préfaisceau $X$, notons $A/X$ sa catégorie des éléments. L'\emph{objet de Lawvere} $L_A$ est le préfaisceau des cribles, c'est-à-dire le classifiant des sous-objets de $\PSh(A)$. Une \emph{catégorie test faible} est une catégorie pour laquelle $A/L_A$ est asphérique, entre autres conditions, et une \emph{catégorie test locale} une catégorie dont les tranches $A/a$ sont des catégories test faibles \cite[déf.~4.1.8, 4.1.12]{cisinski-308}.

\begin{bridge}[Catégories test]\label{br:test}
$A$ est une catégorie test locale si et seulement si $L_A$ est localement asphérique : $A/(a\times L_A)$ est $W$-asphérique pour tout objet $a$ \cite[th.~4.1.19]{cisinski-308}, \cite[th.~1.5.6]{maltsiniotis-301}. Lean : non formalisé ; Mathlib n'a pas de localisateurs fondamentaux. Le fichier \texttt{lean/bridges/Bridges/FundamentalLocalizers.lean} de l'édition démontre les conséquences ci-dessous dans un cadre axiomatique.
\end{bridge}

\begin{theorem}\label{th:testore}
Une petite catégorie non vide qui satisfait la condition d'Ore n'est ni une catégorie test faible ni une catégorie test locale. De façon équivalente, le topos des préfaisceaux sur une catégorie test faible ou locale n'est jamais de De Morgan ; pour une catégorie test locale $A$, aucune tranche $\PSh(A)/X$ avec $X$ non vide n'est de De Morgan non plus.
\end{theorem}

\begin{trace}
\item[tablier] $\PSh(A)$ pour une petite catégorie $A$, avec son classifiant des sous-objets $L_A$.
\item[arche] le pont~\ref{br:dm} d'un côté, le pont~\ref{br:test} de l'autre.
\item[spécialiser] la condition d'Ore dit qu'un crible non vide ne se tire jamais en arrière sur le crible vide ; de façon équivalente, les cribles vides se détachent : $L_A = L_A^{\emptyset}\sqcup L_A^{+}$. Dans le langage des topos, cela dit que $\bot\colon 1\to\Omega$ a un complément, ce qui caractérise les topos de De Morgan \cite[\S1]{caramello-demorgan}, \cite[prop.~14]{caramello-chamoun}.
\item[réécrire] alors $A/(a\times L_A)$ est réunion disjointe de deux morceaux non vides (le crible vide et le crible maximal sur $a$), donc n'est pas connexe ; pour $W$ non grossier, une catégorie asphérique est connexe \cite[lemme~1.7.22]{maltsiniotis-301}.
\item[vérifier] l'argument direct est l'étape de réécriture ; Lean : axiomatique seulement (\lean{not_isDeMorganTopos_slice_of_isLocalTest}, dans \texttt{lean/bridges} de l'édition).
\item[nouveauté] non trouvé tel quel ; voir ci-dessous.
\end{trace}

\noindent Les tranches suivent, puisque $A/X$ est une catégorie test locale dès que $A$ l'est \cite[prop.~4.1.14 (v)]{cisinski-308}, \cite[rem.~1.5.4 b)]{maltsiniotis-301}, et que $\PSh(A)/X\simeq\PSh(A/X)$. Ainsi les ensembles simpliciaux, cubiques, les $\Theta$-ensembles et les ensembles dendroïdaux ne forment pas des topos de De Morgan, ni aucune de leurs tranches ; dans $\Delta$, le cospan $[0]\rightrightarrows[1]$ des deux faces ne se complète pas. La réciproque est fausse : $\Delta$ ne satisfait pas la condition d'Ore et est une catégorie test.

\status{} Non trouvé tel quel, ni dans \cite{cisinski-308}, \cite{maltsiniotis-301}, \emph{Pursuing Stacks}, ni par TheoremSearch et alphaXiv : aucun des textes sur les catégories test ne mentionne la condition d'Ore ni la loi de De Morgan, même si \emph{Pursuing Stacks} \S88 dit sans démonstration qu'une catégorie test n'a pas d'objet initial. L'énoncé est une combinaison immédiate de Caramello--Chamoun (topos de De Morgan si et seulement si $\bot$ a un complément), du théorème~4.1.19 de Cisinski et du lemme~1.7.22 de Maltsiniotis. Le résultat publié le plus proche est celui de Mielke : un topos est de De Morgan si et seulement s'il est \emph{localement homotopiquement trivial}, tout intervalle d'une tranche se scindant en deux morceaux qui séparent ses extrémités, la condition d'Ore en étant la forme préfaisceaux \cite[th.~1.5, cor.~1.6]{mielke}. Ses intervalles sont totalement ordonnés et son homotopie n'est pas l'asphéricité de Grothendieck, de sorte que le théorème~\ref{th:testore} n'est pas un cas du sien. Le mécanisme est pourtant le même : dans un topos de De Morgan, l'objet des valeurs de vérité se scinde en deux morceaux.

\section{Étalonnage: ce que le générateur retrouve}\label{sec:calibration}

Un générateur de théorèmes doit être testé sur des entrées dont les sorties sont connues. Appliqué aux tabliers les plus simples, le nôtre produit les résultats suivants, tous publiés sauf un, que Grothendieck énonce comme exercice dans une note non publiée; trois d'entre eux (le critère pour les monoïdes finis, les points essentiels comme idempotents, les sous-topos des $\mathbb N$-ensembles) nous les avons d'abord crus nouveaux, avant que la recherche bibliographique de l'étape~6 ne les retrouve. Pour les autres lignes, l'article ne consigne pas si nous les avons crus nouveaux. Nous les citons seulement pour montrer que la machine est correcte; nous ne les présentons pas comme nouveaux. L'article ne consigne de l'exécution que ce qui suit.

\begin{center}\small
\begin{tabular}{@{}p{0.8\linewidth}r@{}}
\toprule
élément & nombre\\
\midrule
énoncés du journal d'étalonnage (lignes ci-dessous) & 14\\
énoncés du tableau des résultats du \S\ref{sec:intro} (lignes, certaines regroupant plusieurs énoncés) & 10\\
énoncés d'abord crus nouveaux, puis trouvés publiés & 3\\
couples (tablier, arche) essayés & non consigné\\
\bottomrule
\end{tabular}
\end{center}

\begin{center}\small
\begin{tabularx}{\linewidth}{>{\raggedright\arraybackslash}p{4.3cm}>{\raggedright\arraybackslash}X>{\raggedright\arraybackslash}p{3.3cm}}
\toprule
tablier $\times$ arche & énoncé engendré & publié\\
\midrule
$M$-ensembles $\times$ booléen & tiers exclu si et seulement si $M$ est un groupe (Lean: \lean{isBooleanTopos_singleObj_iff}) & \cite[Th.~2.4]{hemelaer-rogers}\\
$M$-ensembles $\times$ De Morgan & tiers exclu faible si et seulement si $aM\cap bM\neq\emptyset$ (Lean: \lean{isDeMorganTopos_singleObj_iff}) & \cite[Th.~2.31]{hemelaer-rogers}\\
espace d'Alexandrov $\times$ De Morgan, topologie & un préordre est confluent si et seulement si son espace d'Alexandrov est extrêmement discontinu (Lean: \lean{extremallyDisconnected_iff_confluent}) & \cite[Prop.~7.10]{caramello-stone}\\
préordres $\times$ Morita & topos équivalents si et seulement si les réflexions ordonnées sont isomorphes (Lean: \lean{moritaEquivalentCat_preorder_iff}) & \cite[A1.1.9--10]{elephant}, \cite[\S4.1]{caramello-stone}\\
$M$-ensembles finis $\times$ De Morgan & tiers exclu faible si et seulement s'il y a un seul idéal à droite minimal & dual de \cite[Th.~5.21]{gould} avec \cite[Th.~2.31]{hemelaer-rogers}\\
$M$-ensembles $\times$ De Morgan & monoïdes ayant un élément absorbant à gauche ($mz = z$), par ex.\ monoïdes de transitions des automates synchronisants & \cite[Th.~2.2.56]{rogers-thesis}\\
$M$-ensembles $\times$ points & points essentiels $\leftrightarrow$ idempotents à $\mathcal D$ près ($\mathcal J$ pour $M$ fini) & \cite[Cor.~4.2]{rogers-1905}, \cite[Th.~5.6.1]{pirashvili-jncg}\\
$M$-ensembles $\times$ sous-topos & $M$ fini: idéaux bilatères idempotents & \cite[Prop.~6.1.2]{pirashvili-jncg}\\
$\mathbb N$-ensembles $\times$ sous-topos & exactement trois: tout, les $\mathbb Z$-ensembles, le dégénéré & d'après \cite[Cor.~4.3.11]{pirashvili-jpaa}\\
$M$-ensembles $\times$ Morita & $M\sim N$ si et seulement si $N\cong eMe$ avec $1\in MeM$ & \cite{knauer,banaschewski}, \cite[Cor.~6.6]{rogers-1905}\\
théories de Horn $\times$ De Morgan & topos classifiant de De Morgan & \cite[Th.~3.5]{caramello-demorgan}\\
torseurs $\times$ Morita & $G$-torseurs et $H$-torseurs Morita-équivalents si et seulement si $G\cong H$ & classique (le seul point essentiel des $G$-ensembles a pour groupe d'automorphismes $G$, \S\ref{sec:equivalences})\\
quotients $\times$ sous-topos & intersection des anneaux locaux et des anneaux intègres & \cite[\S18]{caramello-lattices}\\
petites catégories $\times$ Morita & $\widehat C$ est le topos de préfaisceaux d'un préordre si et seulement si $C$ est mince (Lean: \lean{exists_preorder_moritaEquivalentCat_iff_isThin}) & non publié comme tel; exercice d'une note non publiée \cite[p.~26]{grothendieck-161-3} (\S\ref{sec:equivalences})\\
\bottomrule
\end{tabularx}
\end{center}

Un résultat mérite un mot, parce que sa lecture est nouvelle alors que les mathématiques sont connues. Pour un système de réécriture abstrait $(S,\to)$, on munit $S$ de la topologie dont les ouverts sont les ensembles de termes stables par réécriture. Alors \emph{le système est confluent (Church--Rosser) si et seulement si cet espace est extrêmement discontinu}: c'est la troisième ligne du tableau appliquée au préordre $\to^*$. La confluence de la $\beta$-réduction dit donc que l'espace des $\lambda$-termes, muni des ensembles stables par $\beta$-réduction comme ouverts, est extrêmement discontinu. Nous n'avons pas trouvé cette lecture dans la littérature sur la réécriture, mais, comme c'est une simple reformulation d'un théorème publié, nous ne la comptons pas comme nouvelle.

\section{Limites de la démarche}\label{sec:limits}

Cette exécution montre que le pipeline fonctionne, mais aussi par où il pèche. Nous passons ses limites en revue dans l'ordre des étapes de la figure~\ref{fig:pipeline}.

\paragraph{Les arches sont classiques.} Les cinq arches de \lp{} dans lean-pool portent sur des topos de \emph{préfaisceaux}: ce sont des résultats classiques de Caramello, Johnstone, Rogers et d'autres, qui ramènent un invariant à une condition sur une petite catégorie. La section~\ref{sec:analysis} emploie deux arches de plus, les caractérisations classiques des topos de faisceaux booléens et de De Morgan d'un espace; elles y sont énoncées, mais non formalisées. Des conditions sur une petite catégorie sont faciles à énoncer et à vérifier, et la plupart de nos résultats ont une preuve directe courte. Font exception le carquois cyclique (une page), $(\mathbb Z/n,\cdot)$ (la classification de Pirashvili et une factorisation de type Ramsey) et l'énoncé sur les compacts séparés (Stone et Nakano), dont les preuves utilisent des théorèmes cités. La profondeur d'un théorème engendré est bornée par celle de son arche et de son tablier \cite[transparent 8]{caramello-mines}; pour obtenir des résultats plus profonds, il faut des arches sur de vrais sites (par exemple les caractérisations d'invariants géométriques sur les sites dues à Caramello \cite{caramello-tst}) et des tabliers qui relient des théories équivalentes de façon non triviale. Notre registre ne contient ni arche sur de vrais sites ni tablier du second genre; la note associée \cite{bridges-notes} en traite un, le topos des pseudo-corps tiré des notes de Grothendieck.

\paragraph{Les tabliers sont choisis à la main.} Les résultats nouveaux viennent de tabliers qui n'avaient pas encore été associés à ces arches: les représentations ensemblistes de carquois et les monoïdes de transformations classiques. Les choisir était l'étape créative, et rien dans le pipeline ne l'automatise. Un registre de la forme \q{la structure $X$ est un préfaisceau sur la forme $C$} est facile à remplir pour les topos de préfaisceaux; pour les topos classifiants de théories, le remplir systématiquement demande de démontrer des équivalences de Morita, ce qui relève de la recherche.

\paragraph{Spécialiser peut être difficile.} Calculer la condition sur la forme était facile pour Ore sur $\B M$ ou sur une catégorie libre, mais a déjà demandé une page pour les topologies sur le carquois cyclique (théorème~\ref{th:cycle}); les topologies sur $\B(\mathbb N^2)$, ou la condition d'Ore sur une catégorie présentée par générateurs et relations, sont de vrais problèmes. Le générateur déplace la difficulté sans la supprimer.

\paragraph{Déplier et réécrire demandent un humain.} La sortie brute est \q{$\mathrm{Paths}(R^{\op})$ vérifie \texttt{RightOreCondition}}; l'énoncé que veut un mathématicien est \q{chaque sommet a au plus une flèche sortante}. C'est cette étape que l'exposé confie aux systèmes de réécriture et aux modèles de langage \cite[transparent 17]{caramello-mines}, et c'est elle qui nous a pris la plus grande partie de notre temps.

\paragraph{Une exécution automatique des premières étapes.} Pour éprouver la tactique que propose l'exposé sur les deux bases de données \cite[transparent 17]{caramello-mines}, nous avons mis les arches booléenne et de De Morgan, et leurs lectures sur les tabliers \texttt{SingleObj M} et \texttt{Paths Q}, dans un jeu de règles Aesop \cite{aesop}, comme règles de normalisation (\lean{lean/bridges/Bridges/Aesop/}). Devant les énoncés déjà démontrés à la main, \texttt{aesop} trouve la preuve de ceux-ci, et de ceux-ci seulement:
\begin{itemize}
\item \emph{joindre et spécialiser} (étapes~1 et~2): les quatre lectures des deux arches sur les deux tabliers, en un appel chacune;
\item \emph{conclure par un fait connu}: le topos booléen d'un groupe, à condition de donner \texttt{Group.isUnit} comme règle;
\item \emph{un transfert de Morita} (un groupoïde reste booléen le long d'une équivalence de topos de préfaisceaux), à condition d'énoncer l'arche à l'envers dans un lemme à part, de déplier l'équivalence et de donner l'invariance de la propriété comme règle vers l'avant.
\end{itemize}
Il ne trouve pas les corollaires dont la condition sur le tablier demande un témoin: De Morgan pour les monoïdes commutatifs ($c=b$, $d=a$) et pour les monoïdes avec zéro ($c=d=0$), ni l'échec de De Morgan pour le monoïde libre sur deux lettres. Chaque échec est consigné dans le fichier sous \texttt{fail\_if\_success}, de sorte que la compilation casse si la frontière bouge. Deux limites de l'outil apparaissent en chemin: une règle \texttt{simp} ne s'emploie pas de droite à gauche, il faut donc écrire l'arche renversée; et une recherche qui échoue laisse le but ouvert, il faut donc lancer les échecs en recherche terminale. L'automatisation couvre ainsi les étapes qui étaient déjà mécaniques; les témoins, les transferts et le choix des tabliers restent à l'humain. Elle ne propose aucun énoncé: elle démontre ceux qu'on lui donne.

\paragraph{Lean vérifie le corollaire, pas sa lecture.} Une preuve Lean que $\mathrm{Paths}(R^{\op})$ vérifie Ore si et seulement si une condition est satisfaite est une garantie sur cet énoncé, dans les conventions de la bibliothèque (côté de l'action, sens des flèches, \texttt{SingleObj} composant selon $f\gg g = g\cdot f$). Que la phrase informelle de l'article dise la même chose, Lean ne le vérifie pas. Nous avons rencontré deux conventions qui changent le côté d'une action de monoïde, et ce côté compte: les $T_n$-ensembles à droite ($ab = a\circ b$) ne forment pas un topos de De Morgan pour $n\ge 2$, puisque les applications constantes $c_i$ vérifient $c_iT_n = \{c_i\}$, alors que les $T_n$-ensembles à gauche en forment un.

\paragraph{Le filtre de nouveauté est faible.} \q{Non trouvé} ne veut pas dire \q{nouveau}. Le jeu d'appariements de l'étape~6 n'atteint que les énoncés qui ont un jumeau Lean. Notre recherche (annexe~\ref{app:novelty}) a couvert les sources évidentes, mais les énoncés élémentaires relèvent souvent du folklore: on les rencontre en passant ou en exercice, et ils ne sont pas indexés. Le journal d'étalonnage compte 14 énoncés, contre 10 lignes dans le tableau des résultats, et trois énoncés que nous avions crus nouveaux se sont révélés publiés (le critère pour les monoïdes finis, les sous-topos des $\mathbb N$-ensembles, les points essentiels comme classes d'idempotents). Nous n'avons pas consigné le nombre de couples (tablier, arche) essayés (\S\ref{sec:calibration}): ces chiffres ne donnent donc aucune proportion, mais ils montrent que le filtre de nouveauté a changé ce que nous rapportons.

\paragraph{La plupart des sorties sont sans intérêt.} Associer chaque tablier à chaque arche produit beaucoup d'énoncés vrais, nouveaux au sens étroit, et sans aucun intérêt (le tiers exclu faible pour les préfaisceaux sur une catégorie particulière à cinq objets). L'avertissement de l'exposé sur les résultats \q{bizarres} s'applique ici. Pour être utile, un générateur doit classer ses sorties; cela relève du goût mathématique, et ce n'est pas formalisé.

\paragraph{À quoi sert la démarche.} Dans ces limites, le pipeline sert à quelque chose: il pose des questions que personne n'avait posées (le tiers exclu faible est-il vrai pour les automates? combien de sous-topos un cycle d'applications a-t-il?), il y répond avec une garantie (l'arche est démontrée formellement) et il indique une preuve directe. Comme outil d'enseignement, ou pour explorer la frontière entre deux domaines, c'est déjà appréciable.

\subsection{Leçons pour un générateur automatique}

\begin{enumerate}
\item \emph{Choisir les tabliers est l'étape créative.} Les résultats nouveaux viennent de tabliers que nous n'avons pas trouvés associés à ces arches: carquois et monoïdes de transformations. Une base de données de tabliers utile est une base d'énoncés de la forme \q{cette structure familière est un préfaisceau sur telle forme}.
\item \emph{L'étape de dépliage est mécanique.} Le lemme~\ref{lem:neg} est la seule règle nécessaire pour traduire \q{De Morgan} en un énoncé sur des états et des pièges. Un générateur a besoin d'une telle règle par invariant.
\item \emph{C'est la réécriture qui demande le travail.} La sortie brute pour les automates est \q{$\B(\Sigma^*)$ vérifie \texttt{RightOreCondition}}; l'énoncé lisible porte sur les flèches sortantes d'un sommet (théorème~\ref{th:quivers}). C'est l'étape que l'exposé confie aux systèmes de réécriture et aux modèles de langage \cite[transparent 17]{caramello-mines}.
\item \emph{Le filtrage par la littérature est nécessaire.} Ce n'est qu'à ce stade que nous avons trouvé publiés trois de nos candidats (\S\ref{sec:calibration}). Faute d'avoir consigné le nombre de couples essayés, nous ne pouvons donner aucune proportion. Sans filtre de nouveauté, un générateur présente des résultats connus comme nouveaux.
\item \emph{La vérification.} Chaque résultat présenté ici a une preuve directe: courte pour la plupart, d'une page pour le carquois cyclique, plus longue pour $(\mathbb Z/n,\cdot)$ et pour l'énoncé sur les compacts séparés, qui utilisent des théorèmes cités (la classification de Pirashvili, avec une factorisation de type Ramsey; Stone et Nakano). Pour des tabliers plus profonds, une preuve directe peut ne pas exister, et la preuve formelle de l'arche en Lean sera alors la seule garantie. C'est pourquoi on formalise d'abord les arches.
\end{enumerate}

\appendix
\section{Le filtre de nouveauté en détail}\label{app:novelty}

Cette annexe donne en entier l'étape~6 du \S\ref{sec:method}, ainsi que les sources consultées pour les résultats.

\subsection{Le jeu de données d'appariements de TheoremGraph}

On commence par un passage mécanique. On résume en une phrase l'énoncé Lean du théorème réécrit et l'on cherche ce résumé dans le jeu de données d'appariements de TheoremGraph (fiche du jeu de données, Hugging Face, consultée le 23 novembre 2026) \cite{theoremgraph-matching}. Chacune de ses 100\,831 lignes associe une déclaration Lean, de Mathlib ou de l'un de 25 autres projets, à l'énoncé d'arXiv le plus proche, trouvé par similarité cosinus (au moins 0,80) entre plongements de résumés, et classé \emph{exact}, \emph{inexact} ou \emph{faux} par un modèle de langage; son juge principal tient 47\,952 paires pour des appariements. Une réponse positive établit qu'un jumeau formel existe et désigne l'article qui en donne la version informelle; le résultat va alors au journal d'étalonnage, avec une référence précise. Une réponse négative ne prouve rien: un énoncé informel sans déclaration Lean ne figure pas dans le jeu, les étiquettes sont celles d'un modèle et ne sont pas vérifiées, et le texte informel n'est fourni que pour les paires dont la source arXiv autorise la redistribution. La recherche à la main reste donc nécessaire. Nous n'avons fait ce passage que sur huit résultats (les énoncés sur les carquois, le carquois cyclique, Morita, les coins, $T_n$, $M_n(k)$, $(\mathbb Z/n,\cdot)$ et les compacts séparés; ni les automates, ni les topos combinatoires, ni les treillis de sous-espaces), par mots-clés, dans le service de recherche public du jeu, et aucune des paires renvoyées ne concernait l'un d'eux. On s'y attendait, et cela ne dit rien de la nouveauté: le côté Lean du jeu ne contient aucune déclaration de Lean Pool, ni aucune sur les topos de De Morgan, et une recherche par mots-clés classe selon les mots, non selon la similarité des résumés.

\subsection{La recherche sémantique de TheoremSearch, et l'OEIS}

Pour le côté informel, nous avons ensuite soumis les dix lignes du tableau des résultats à la recherche sémantique de TheoremSearch \cite{theoremsearch}, qui plonge une description en langue naturelle et classe quelque 8 millions d'énoncés tirés d'arXiv et d'autres sources selon la similarité de leurs résumés (API publique, 23 novembre 2026, six résultats par requête). Aucun des énoncés renvoyés ne formule l'un des dix. Trois requêtes en ont renvoyé les ingrédients, que l'article signale déjà: le théorème fondamental de la géométrie projective, pour les treillis de sous-espaces; pour le résultat sur les compacts séparés, la caractérisation des topos localiques de De Morgan (celle de Johnstone, sous la forme que lui donne Caramello) et celle des compacts séparés dont $C(X,\mathbb R)$ est complet au sens de Dedekind comme les espaces stoniens, en deux énoncés distincts; pour les coins, des énoncés sur les anneaux où une propriété de $R$ passe à $eRe$ quand $ReR = R$, mais pas l'énoncé sur les monoïdes. Les autres requêtes ont renvoyé des énoncés sur les carquois, les cycles, les monoïdes ou les diviseurs, sans rapport avec nos résultats. Avec six résultats par requête, la recherche reste peu profonde, et un énoncé formulé dans un autre vocabulaire peut lui échapper; ce passage resserre donc la recherche à la main sans la remplacer. Les appels, leurs paramètres et les réponses du 23 novembre 2026 sont dans \texttt{docs/bridges/novelty/} (\texttt{theoremsearch.py} et \texttt{theoremsearch-2026-11-23.json}): une requête \texttt{POST https://api.theoremsearch.com/search} par résultat, avec la description comme \texttt{query}, \texttt{n\_results} fixé à 6, sans filtre ni pondération par les citations. Les résultats numériques ont aussi été cherchés dans l'OEIS \cite{oeis} (23 novembre 2026): les $2^n+1$ sous-topos du $n$-cycle sont A000051; les $2^{\omega(n)}$ points essentiels des $(\mathbb Z/n,\cdot)$-ensembles sont A034444, le nombre de diviseurs unitaires de $n$, comme il se doit puisque les idempotents de $\mathbb Z/n$ correspondent à ses diviseurs unitaires; les nombres de Dedekind sont A000372; la suite $D(\omega(n))$ des sous-topos des $(\mathbb Z/n,\cdot)$-ensembles ($2, 3, 3, 3, 3, 6, 3, \dots$ pour $n=1,\dots,16$) n'est pas dans l'OEIS, bien qu'elle soit la composée de deux suites qui y sont.

\subsection{La recherche à la main}

On cherche ensuite, dans cet ordre: l'\emph{Elephant} de Johnstone; les articles de l'auteur de l'arche (ici Caramello); la littérature récente sur le tablier (ici Hemelaer, Rogers, Pirashvili pour les actions de monoïdes); la littérature classique du domaine de l'énoncé réécrit (demi-groupes, carquois, topologie); les serveurs de prépublications et le nLab. Il faut chercher l'énoncé \emph{réécrit}, dans le vocabulaire de son domaine, et pas seulement l'énoncé topos-théorique: trois de nos candidats (le critère pour les monoïdes finis, les sous-topos des $\mathbb N$-ensembles, les points essentiels comme classes d'idempotents) étaient publiés dans le langage des demi-groupes ou des actions de monoïdes, alors que nous les avions crus nouveaux. Les énoncés trouvés sont consignés au journal d'étalonnage (\S\ref{sec:calibration}).

\subsection{Sources consultées pour les résultats}

Pour chaque énoncé des \S\ref{sec:machines}--\ref{sec:shapes}, nous avons consulté l'\emph{Elephant} de Johnstone, les articles et le livre de Caramello (dont \cite{caramello-chamoun}), les articles et la thèse de Hemelaer et Rogers sur les topos d'actions de monoïdes \cite{hemelaer-rogers,rogers-1905,rogers-thesis}, la littérature sur l'équivalence de Morita et les relations de Green des monoïdes, la littérature sur les automates synchronisants, le nLab, ainsi que des recherches générales sur le web et dans les serveurs de prépublications. \q{Non trouvé} signifie exactement cela; plusieurs énoncés sont élémentaires et relèvent peut-être du folklore, et nous accueillerons volontiers des références.

\section{Lean}\label{sec:lean-lit}

Les cinq arches de \lp{}, dans lean-pool, sont démontrées en Lean~4 avec Mathlib, sans autre axiome que \texttt{propext}, \texttt{Classical.choice} et \texttt{Quot.sound}. Cette affirmation décrit le développement: les fichiers ne contiennent aucune commande \texttt{\#print axioms}, et nous ne consignons aucune vérification séparée à leur sujet. Pour \texttt{lean/bridges}, l'intégration continue du dépôt de cette édition fait un audit des axiomes à chaque modification: à la version décrite ici, il portait sur 437 déclarations, toutes limitées à \texttt{propext}, \texttt{Classical.choice} et \texttt{Quot.sound}. Les énoncés engendrés sont formalisés, comme corollaires des arches, dans le paquet Lake \texttt{lean/bridges} du dépôt de cette édition (\url{https://github.com/Commutator-IO/grothendieck-project}), qui dépend de lean-pool à une révision fixée. Les fichiers utilisés ici sont les suivants; le paquet en contient trois autres, décrits dans la note associée \cite{bridges-notes}:
\begin{itemize}
\item \lean{Monoids.lean} (\lean{isBooleanTopos_singleObj_iff}, \lean{isDeMorganTopos_singleObj_iff}, \lean{isDeMorganTopos_singleObj_of_monoidWithZero}, \lean{isDeMorganTopos_singleObj_of_commMonoid}, \lean{isDeMorganTopos_singleObj_freeMonoid_unit}, \lean{not_isDeMorganTopos_freeMonoid_two}, \lean{isBooleanTopos_singleObj_of_group}); \lean{Preorders.lean} (\lean{isDeMorganTopos_functor_preorder_iff}, \lean{extremallyDisconnected_iff_confluent}, \lean{moritaEquivalentCat_preorder_iff}); \lean{Combinatorics.lean} (\lean{isDeMorganTopos_presheaf_of_hasInitial}, \lean{not_isDeMorganTopos_graphs}, \lean{graphs_pseudocompl_counterexample}, \lean{not_isDeMorganTopos_simplicialSets}, \lean{isDeMorganTopos_augmentedSimplicialSets}, \lean{isDeMorganTopos_presheaf_fintypeCatSkeleton}, \lean{isDeMorganTopos_treeTopos}, \lean{not_isBooleanTopos_treeTopos}); \lean{EssentialPoints.lean} (\lean{nonempty_essentialPoint_iso_iff}, \lean{exists_iso_essentialPoint}, \lean{essentialPointEndMulEquivOfGroup}, \lean{karoubiRel_iff_rank_eq}, \lean{essentialPointClassesEquivFin}).
\item \lean{QuiversMorita.lean} (Morita et points essentiels des représentations de carquois: \lean{moritaEquivalentCat_paths_iff}, \lean{nonempty_representations_equiv_iff}, \lean{essentialPointPathsEquiv}) et \lean{QuiversCyclic.lean} (les $2^n+1$ topologies du carquois cyclique: \lean{card_grothendieckTopology}).
\item \lean{QuiversDeMorgan.lean} (De Morgan et booléen pour les représentations de carquois: \lean{rightOreCondition_paths_iff}, \lean{isDeMorganTopos_paths_iff}, \lean{isBooleanTopos_paths_iff}; il ne contient aucun énoncé de Morita ni de points essentiels), \lean{Corners.lean} (transfert aux coins), \lean{KaroubiMonoid.lean} (les $n$ classes pour $T_n$) et, pour le journal d'étalonnage, \lean{FiniteMonoids.lean} (unique idéal à droite minimal) et \lean{NatTopologies.lean} (les trois topologies sur $\B\mathbb N$, \lean{nat_card_grothendieckTopology_natCat}, avec les faisceaux de celle du milieu, \lean{isSheaf_finalTopology_iff}).
\end{itemize}
Ne sont pas formalisés: le chapitre d'analyse (Mathlib n'a pas le topos des faisceaux sur un espace avec sa négation), les treillis de sous-espaces au-delà des énoncés de De Morgan et de Morita, les dénombrements pour $PT_n$, $I_n$, $O_n$ et $(\mathbb Z/n,\cdot)$, tous les dénombrements de sous-topos de monoïdes, et la description des faisceaux du carquois cyclique. Les énumérations par ordinateur pour $(\mathbb Z/n,\cdot)$ (\S\ref{sec:linear}) ne font pas partie du développement Lean; pour $n>30$, elles passent par la classification des idéaux idempotents et ne vérifient donc pas les topologies indépendamment.
Dans Mathlib, \texttt{SingleObj M} compose selon $f\gg g = g\cdot f$; avec cette convention, la condition d'Ore sur \texttt{SingleObj M} s'écrit $\forall a\,b,\ \exists c\,d,\ a\cdot c = b\cdot d$ (\lean{rightOreCondition_singleObj_iff}), comme dans le texte.

\begin{center}\scriptsize\emergencystretch=3em
\begingroup\def\UrlBreaks{\do\_\do\.\do\/}
\begin{tabularx}{\linewidth}{>{\raggedright\arraybackslash}p{2.6cm}>{\raggedright\arraybackslash}X>{\raggedright\arraybackslash}p{2.9cm}>{\raggedright\arraybackslash}p{2.2cm}}
\toprule
énoncé & déclaration Lean & fichier & statut\\
\midrule
$M$-ensembles booléens si et seulement si groupe; De Morgan si et seulement si Ore & \lean{isBooleanTopos_singleObj_iff}, \lean{isDeMorganTopos_singleObj_iff} & \lean{Monoids} & démontré en Lean\\
carquois: De Morgan, booléen (th.~\ref{th:quivers}) & \lean{isDeMorganTopos_paths_iff}, \lean{isBooleanTopos_paths_iff} & \lean{QuiversDeMorgan} & démontré en Lean (convention de Lean)\\
automates sur $\Sigma$ & \lean{isDeMorganTopos_singleObj_freeMonoid_unit}, \lean{not_isDeMorganTopos_freeMonoid_two} & \lean{Monoids} & démontré en Lean\\
carquois: Morita, points essentiels & \lean{nonempty_representations_equiv_iff}, \lean{essentialPointPathsEquiv} & \lean{QuiversMorita} & démontré en Lean\\
carquois cyclique: $2^n+1$ topologies & \lean{card_grothendieckTopology} & \lean{QuiversCyclic} & démontré en Lean pour le nombre de topologies; les faisceaux de chacune, à la main seulement\\
$T_n$: $n$ points essentiels & \lean{karoubiClassesEquivFin} & \lean{KaroubiMonoid} & démontré en Lean (classes de l'enveloppe de Karoubi)\\
$M_n(\mathbb F_q)$: $n+1$ points essentiels & \lean{essentialPointClassesEquivFin} & \lean{EssentialPoints} & démontré en Lean, sur un corps quelconque\\
points essentiels de $PT_n$, $I_n$, $O_n$; tous les dénombrements de sous-topos du \S\ref{sec:finite} & --- & --- & à la main seulement\\
\bottomrule
\end{tabularx}

\begin{tabularx}{\linewidth}{>{\raggedright\arraybackslash}p{2.6cm}>{\raggedright\arraybackslash}X>{\raggedright\arraybackslash}p{2.9cm}>{\raggedright\arraybackslash}p{2.2cm}}
\toprule
énoncé & déclaration Lean & fichier & statut\\
\midrule
coins: Ore se transfère & \lean{forall_exists_mul_eq_mul_iff_cornerMonoid} & \lean{Corners} & démontré en Lean\\
$M_n(k)$-ensembles, $\mathsf{Vect}_{\le n}$ & --- & --- & à la main seulement (les classes de rang sont en Lean, voir plus haut)\\
treillis de sous-espaces: De Morgan, Morita & \lean{isDeMorganTopos_functor_preorder_iff}, \lean{moritaEquivalentCat_preorder_iff} & \lean{Preorders} & démontré en Lean; le reste à la main seulement\\
dénombrements pour $(\mathbb Z/n,\cdot)$ & --- & --- & énumérés par ordinateur (pas en Lean); preuve à la main seulement\\
$X$ compact séparé & --- & --- & à la main seulement\\
graphes, ensembles simpliciaux, augmentés, $\mathsf{Fin}$, arbres & \lean{not_isDeMorganTopos_graphs}, \lean{not_isDeMorganTopos_simplicialSets}, \lean{isDeMorganTopos_augmentedSimplicialSets}, \lean{isDeMorganTopos_presheaf_fintypeCatSkeleton}, \lean{isDeMorganTopos_treeTopos} & \lean{Combinatorics} & démontré en Lean; les autres lignes combinatoires à la main seulement\\
\bottomrule
\end{tabularx}
\endgroup
\end{center}

% Bibliographie commune, en anglais : composée en anglais pour éviter les espaces françaises avant « : » et « ; ».
\begin{english}\renewcommand{\refname}{Références}
\begin{thebibliography}{99}
\bibitem{caramello-mines} O.~Caramello, \emph{Topos-empowered automatic theorem generation}, talk, CRI, Mines Paris, 17 September 2024. \url{https://www.oliviacaramello.com/Talks/MinesTalk.pdf}
\bibitem{caramello-unification} O.~Caramello, \emph{The unification of Mathematics via Topos Theory}, arXiv:1006.3930, 2010.
\bibitem{caramello-tst} O.~Caramello, \emph{Theories, Sites, Toposes: Relating and studying mathematical theories through topos-theoretic `bridges'}, Oxford University Press, 2017.
\bibitem{caramello-demorgan} O.~Caramello, \emph{De Morgan classifying toposes}, Advances in Mathematics 222 (2009) 2117--2144; arXiv:0808.1519.
\bibitem{elephant} P.~T.~Johnstone, \emph{Sketches of an Elephant: A Topos Theory Compendium}, Oxford University Press, 2002.
\bibitem{johnstone-demorgan} P.~T.~Johnstone, \emph{Conditions related to De Morgan's law}, in \emph{Applications of Sheaves}, Lecture Notes in Math.~753, Springer, 1979.
\bibitem{hemelaer-rogers} J.~Hemelaer, M.~Rogers, \emph{Monoid properties as invariants of toposes of monoid actions}, Applied Categorical Structures (2021), doi:10.1007/s10485-020-09620-y; arXiv:2004.10513.
\bibitem{rogers-thesis} M.~Rogers, \emph{Toposes of monoid actions}, PhD thesis, University of Insubria, 2021; arXiv:2112.10198.
\bibitem{knauer} U.~Knauer, \emph{Projectivity of acts and Morita equivalence of monoids}, Semigroup Forum 3 (1972).
\bibitem{banaschewski} B.~Banaschewski, \emph{Functors into categories of $M$-sets}, Abh. Math. Sem. Univ. Hamburg 38 (1972).
\bibitem{rogers-1905} M.~Rogers, \emph{Toposes of monoid actions}, arXiv:1905.10277, 2019.
\bibitem{costa-steinberg} A.~Costa, B.~Steinberg, \emph{The Schützenberger category of a semigroup}, arXiv:1408.1615, 2014.
\bibitem{caramello-stone} O.~Caramello, \emph{A topos-theoretic approach to Stone-type dualities}, arXiv:1103.3493, 2011.
\bibitem{caramello-lattices} O.~Caramello, \emph{Lattices of theories}, arXiv:0905.0299, 2009.
\bibitem{caramello-chamoun} O.~Caramello, R.~Chamoun, \emph{De Morgan's law in toposes I}, arXiv:2507.01869, 2025.
\bibitem{chagrov} A.~Chagrov, M.~Zakharyaschev, \emph{Modal Logic}, Oxford University Press, 1997.
\bibitem{clifford-preston} A.~H.~Clifford, G.~B.~Preston, \emph{The Algebraic Theory of Semigroups}, vol.~I, AMS, 1961.
\bibitem{ganyushkin} O.~Ganyushkin, V.~Mazorchuk, \emph{Classical Finite Transformation Semigroups}, Springer, 2009.
\bibitem{barendregt} H.~P.~Barendregt, \emph{The Lambda Calculus: Its Syntax and Semantics}, North-Holland, 1984.
\bibitem{birkedal} L.~Birkedal, R.~E.~Møgelberg, J.~Schwinghammer, K.~Støvring, \emph{First steps in synthetic guarded domain theory: step-indexing in the topos of trees}, Logical Methods in Computer Science 8 (2012).
\bibitem{gould} V.~Gould, C.~Miller, T.~Quinn-Gregson, N.~Ruškuc, \emph{On minimal ideals in pseudo-finite semigroups}, Canad. J. Math. (2023); arXiv:2204.10155.
\bibitem{pirashvili-jpaa} I.~Pirashvili, \emph{Reconstruction theorem for monoid schemes}, J. Pure Appl. Algebra 225 (2021) 106782; arXiv:2002.04336.
\bibitem{pirashvili-jncg} I.~Pirashvili, \emph{The points and localisations of the topos of $M$-sets}, J. Noncommut. Geom. 18 (2024); arXiv:2011.11747.
\bibitem{sga1} A.~Grothendieck, \emph{Revêtements étales et groupe fondamental} (SGA~1), Lecture Notes in Math.~224, Springer, 1971, exposé~V.
\bibitem{szamuely} T.~Szamuely, \emph{Galois Groups and Fundamental Groups}, Cambridge University Press, 2009.
\bibitem{schmid} Schmid, Tohmé, Chin, \emph{Sheaf subcategories of quiver representations}, arXiv:2510.23580, version~2, 2026.
\end{thebibliography}

\end{english}
\end{document}
